% From pandas (and Polars) to simplydata: written by hand, read by docs/make_manual.py (% From pandas (and Polars) to simplydata: written by hand, read by docs/make_manual.py (% From pandas (and Polars) to simplydata: written by hand, read by docs/make_manual.py (% From pandas (and Polars) to simplydata: written by hand, read by docs/make_manual.py (\input{from_pandas}).
% The topics follow the order of the pandas User Guide and the Polars User Guide, used as a checklist only: the text
% and the examples are our own, and each section links to the official page. Every example in a py box is run by
% tests/test_from_pandas_chapter.py (with "import simplydata as sdata" in front), so it cannot go out of date.
% A line "% needs: duckdb" just before a box says the box needs that package (pip install "simplydata[sql]").

\unchapter{From pandas (and Polars) to simplydata}{from-pandas}

This chapter is for you if you already know pandas or Polars, or if you want to learn them after
simplydata. It walks through the topics of the two official guides, in their order:
the \href{https://pandas.pydata.org/docs/user_guide/index.html}{pandas User Guide} and the
\href{https://docs.pola.rs/user-guide/getting-started/}{Polars User Guide}. For each topic it says
what the job is, which simplydata words do it, and a short example on the practice tables. Some jobs
are left to pandas and Polars on purpose. For those, the chapter shows how to hand the table over and
how to come back.

Every example starts with \code{import simplydata as sdata}. It is written once here, not in each
box. The examples use the practice tables of \code{sdata.get\_example}, so you can run them as they
are.

\unsection{Five things to know first}{from-pandas:first}

\begin{enumerate}
  \item \textbf{A word changes the table.} In pandas and Polars a method gives back a new table, so you
    write \code{df = df.dropna()}. In simplydata, \code{table.remove\_rows\_with\_empty\_cells\_in()}
    changes \code{table} itself, and \code{table.undo()} takes the step back. When you want to keep the
    old table too, make a copy first with \code{table.make\_copy()}. Words that find an answer
    (\code{find\_average\_of}, \code{find\_summary\_of}, \code{run\_sql} and the other
    \code{find\_} words) leave the table as it is: they say the answer and give back a new table with it,
    as in \code{averages = table.find\_average\_of("marks", for\_each\_group\_in="class")}.
  \item \textbf{Rows are numbered from 1.} pandas \code{iloc} and Polars \code{row} count from 0.
    In simplydata the first row is row 1, as people count.
  \item \textbf{There is no row index.} A table is only its columns. Polars works the same way. Where
    pandas uses the index (to select, to sort, to line up two tables), simplydata uses a column.
  \item \textbf{Each word says what it did.} Most words print a sentence: how many rows were kept,
    which rows did not match in a join, what the average is. Read it; it often answers the next
    question.
  \item \textbf{The door is always open.} \code{table.as\_pandas()} gives a pandas DataFrame and
    \code{table.as\_polars()} a Polars DataFrame. \code{sdata.table(df)} turns either one back into a
    simplydata table. Use this for any job this chapter leaves to pandas or Polars.
\end{enumerate}

\begin{lstlisting}[style=py]
table = sdata.get_example("marks")
df = table.as_pandas()                   # a pandas DataFrame: use any pandas code here
df["marks_out_of_10"] = df["marks"] / 10
table = sdata.table(df)                  # back to simplydata
table.show_first_rows(3)
\end{lstlisting}

Do you know a pandas or Polars name, but not the simplydata word? Ask: \code{sdata.help("groupby")}.
Names inside the text and date steps work too: \code{sdata.help("str.replace")},
\code{sdata.help("dt.year")} and \code{sdata.help("groupby.filter")} each name the simplydata word.
The full list is in the section ``Coming from another tool?'' of the simplydata dictionary
(\href{https://github.com/bioforkgithub/simplydata/blob/main/docs/manual/simplydataDictionary.pdf}%
{simplydataDictionary.pdf}).

\unsection{A first look at a table}{from-pandas:first-look}

\textbf{The job.} Open a table, look at the first rows, see its size, the kinds of its columns and a
short summary. In the guides: \href{https://pandas.pydata.org/docs/user_guide/10min.html}{10 minutes
to pandas}, \href{https://pandas.pydata.org/docs/user_guide/dsintro.html}{Intro to data structures},
\href{https://pandas.pydata.org/docs/user_guide/basics.html}{Essential basic functionality};
\href{https://docs.pola.rs/user-guide/getting-started/}{Getting started} and
\href{https://docs.pola.rs/user-guide/concepts/data-types-and-structures/}{Data types and structures}
in Polars.

\textbf{In simplydata.} \code{show\_first\_rows} is pandas \code{head}; \code{count\_rows} and
\code{count\_columns} are \code{shape}; \code{describe\_columns} is a mix of \code{info} and
\code{describe}; \code{find\_summary\_of} gives the measures you choose.

\begin{lstlisting}[style=py]
table = sdata.get_example("marks")
table.show_first_rows(5)
table.count_rows()
table.count_columns()
table.get_column_names()
table.describe_columns()
table.find_summary_of("marks", show=["average", "median", "standard deviation"])
\end{lstlisting}

A pandas Series is one column. In simplydata, \code{table.get\_column("marks")} gives the column as a
list of numbers with its own words (\code{find\_average}, \code{find\_percentile} and more), or as a
plain Python list when the column holds text.

\begin{lstlisting}[style=py]
table = sdata.get_example("marks")
marks = table.get_column("marks")
marks.find_average()
names = table.get_column("name")         # text: a plain Python list
print(names[:3])
\end{lstlisting}

\unsection{Kinds of columns}{from-pandas:kinds}

\textbf{The job.} Make a column hold numbers, text, dates or categories. pandas calls this
\code{astype} and \code{to\_numeric}; Polars calls it \code{cast}
(\href{https://docs.pola.rs/user-guide/expressions/casting/}{Casting}).

\textbf{In simplydata.} One word for each kind: \code{turn\_into\_numbers},
\code{turn\_into\_text}, \code{turn\_into\_dates} and \code{turn\_into\_categories}. Each says how
many cells could not be turned and became empty, so nothing is lost without a word.

\begin{lstlisting}[style=py]
table = sdata.get_example("shop")
table.turn_into_dates("date")
table.turn_into_text("customers")
table.turn_into_numbers("customers")
table.describe_columns()
\end{lstlisting}

\textbf{Left to pandas and Polars.} simplydata chooses the storage type for you (64-bit numbers,
pandas' nullable types, Arrow text). Two jobs about storage are left out on purpose:
\begin{itemize}
  \item \textbf{Smaller types to save memory} (pandas \code{to\_numeric(downcast=)}, Polars
    \code{shrink\_dtype}). A research table rarely needs it.
  \item \textbf{Memory use} (pandas \code{memory\_usage(deep=True)}, Polars \code{estimated\_size}).
\end{itemize}
Hand the table over when you need them:

\begin{lstlisting}[style=py]
import polars as pl

table = sdata.get_example("river")
df = table.as_polars()
print(df.estimated_size("kb"), "kB")
small = pl.DataFrame([df[name].shrink_dtype() if df[name].dtype.is_numeric() else df[name]
                      for name in df.columns])
print(small.estimated_size("kb"), "kB")
table = sdata.table(small)
\end{lstlisting}

\unsection{Reading and saving files}{from-pandas:files}

\textbf{The job.} Read a file into a table and save a table to a file. pandas has one function for each
file type (\href{https://pandas.pydata.org/docs/user_guide/io.html}{IO tools}); so has Polars
(\href{https://docs.pola.rs/user-guide/io/}{IO}).

\textbf{In simplydata.} One word each way. \code{sdata.table("file")} reads .csv, .tsv, .xlsx,
.parquet, .json, .jsonl, .feather, SPSS (.sav), Stata (.dta), database files and more: the end of the
name says the type. \code{table.save("file")} writes them. \code{sdata.run\_sql} and
\code{table.run\_sql} run SQL with DuckDB.

\begin{lstlisting}[style=py]
table = sdata.get_example("marks")
table.save("marks.csv")
table.save("marks.xlsx", in_sheet="2026")
table.save("marks.parquet")
again = sdata.table("marks.parquet")
again.count_rows()
\end{lstlisting}

% needs: duckdb
\begin{lstlisting}[style=py]
river = sdata.get_example("river")
averages = river.run_sql("SELECT station, avg(nitrate) AS nitrate FROM data GROUP BY station ORDER BY station")
averages.show()
\end{lstlisting}

\textbf{A table for a report (Markdown or LaTeX).} simplydata saves .html tables, but not Markdown
(.md) or LaTeX (.tex). pandas has \code{to\_markdown}, which needs the package tabulate, and
\code{to\_latex}, which needs jinja2. You can make both without extra packages: Polars prints a
Markdown table, and a few lines of Python write a LaTeX one.

\begin{lstlisting}[style=py]
import polars as pl

table = sdata.get_example("marks")
averages = table.find_average_of("marks", for_each_group_in="class")    # one row for each class
df = averages.as_polars()
with pl.Config(tbl_formatting="MARKDOWN", tbl_hide_column_data_types=True,
               tbl_hide_dataframe_shape=True, tbl_rows=-1):
    markdown = str(df)
with open("averages.md", "w", encoding="utf-8") as file:
    file.write(markdown + "\n")
print(markdown)
\end{lstlisting}

\begin{lstlisting}[style=py]
table = sdata.get_example("marks")
averages = table.find_average_of("marks", for_each_group_in="class")
columns = averages.get_column_names()
lines = [r"\begin{tabular}{" + "l" * len(columns) + "}", " & ".join(columns) + r" \\ \hline"]
for row in averages.as_list():            # one dict for each row
    lines.append(" & ".join(str(row[name]) for name in columns) + r" \\")
lines.append(r"\end{tabular}")
with open("averages.tex", "w", encoding="utf-8") as file:
    file.write("\n".join(lines) + "\n")
\end{lstlisting}

\textbf{Left to pandas and Polars.} These reading jobs are rare in research, so the manual covers them
here instead of a word:
\begin{itemize}
  \item \textbf{Skip broken lines} (pandas \code{on\_bad\_lines="skip"}). simplydata stops and names the
    broken line instead, because skipping loses data without a word. If you must skip, read with pandas
    and give the result to \code{sdata.table}.
  \item \textbf{Nested JSON} (pandas \code{json\_normalize}, Polars \code{unnest}).
  \item \textbf{The tables of a web page} (pandas \code{read\_html}). simplydata reads files at a web
    address, but not the tables inside a page.
  \item \textbf{The clipboard} (pandas \code{read\_clipboard}).
  \item \textbf{A file too big for memory}, read in pieces or lazily (pandas \code{chunksize=}, Polars
    \code{scan\_csv} and the \href{https://docs.pola.rs/user-guide/concepts/lazy-api/}{Lazy API}).
    simplydata works on tables that fit in memory; \code{fast=True} makes it use Polars underneath.
\end{itemize}

\begin{lstlisting}[style=py]
import pandas as pd

records = [{"station": "St.1", "depth": {"top": 1, "bottom": 10}},
           {"station": "St.2", "depth": {"top": 2, "bottom": 25}}]
flat = pd.json_normalize(records)        # columns station, depth.top, depth.bottom
table = sdata.table(flat)
table.rename_columns({"depth.top": "top", "depth.bottom": "bottom"})
table.show()
\end{lstlisting}

\unsection{Choosing rows and columns}{from-pandas:choosing}

\textbf{The job.} Keep some rows, some columns, or one cell. In pandas this is
\href{https://pandas.pydata.org/docs/user_guide/indexing.html}{Indexing and selecting data}
(\code{loc}, \code{iloc}, \code{query}); in Polars it is \code{filter} and \code{select}
(\href{https://docs.pola.rs/user-guide/expressions/basic-operations/}{Basic operations} and
\href{https://docs.pola.rs/user-guide/expressions/expression-expansion/}{Expression expansion}).

\textbf{In simplydata.} \code{keep\_rows\_where} keeps rows by a condition (or a written condition in
quotes); \code{keep\_columns} keeps columns by name, by kind or by a part of the name;
\code{keep\_rows\_by\_position} is \code{iloc} with rows from 1; \code{get\_cell\_at} is one cell.

\begin{lstlisting}[style=py]
table = sdata.get_example("marks")
table.keep_rows_where("subject", is_equal_to="Maths")
table.keep_rows_where("marks > 60 and class == 'A'")
table.keep_columns("name", "class", "marks")
table.show()
table.get_cell_at(row=1, column="marks")
\end{lstlisting}

\begin{lstlisting}[style=py]
weather = sdata.get_example("weather")
weather.keep_columns(starts_with="wind")      # pandas df.filter(regex="^wind"), Polars pl.col("^wind.*$")
weather.keep_rows_by_position(1, to=3)        # pandas iloc[0:3]
weather.show()
\end{lstlisting}

\textbf{Left to pandas and Polars: the row index.} simplydata has no row index, so these jobs are left
out on purpose. Each has a route with a column.
\begin{itemize}
  \item \textbf{Set or reset an index, choose rows by label, slice labels} (pandas \code{set\_index},
    \code{reset\_index}, \code{loc["a":"c"]},
    \href{https://pandas.pydata.org/docs/user_guide/advanced.html}{MultiIndex}).
    Route: keep the labels in a column and use \code{keep\_rows\_where}.
  \item \textbf{Sort by the index} (pandas \code{sort\_index}). Route: \code{sort\_by} the column.
  \item \textbf{Repeated labels} (pandas \href{https://pandas.pydata.org/docs/user_guide/duplicates.html}{Duplicate
    Labels}). Route: \code{show\_duplicate\_rows(columns=[...])} on the column.
  \item \textbf{Choose rows with a function inside a chain} (pandas \code{loc[lambda d: ...]}). Each
    simplydata word is its own line, so the table you pass is always the table as it is now; a written
    condition, \code{keep\_rows\_where("marks > hours * 10")}, does the same job.
\end{itemize}
When you need the index itself, hand the table over:

\begin{lstlisting}[style=py]
table = sdata.get_example("weather")
df = table.as_pandas().set_index(["city", "month"])     # a MultiIndex
cold = df.loc["Frostby"].sort_index()                   # rows by label, sorted by the index
table = sdata.table(cold.reset_index())                 # the index becomes a column again
table.show_first_rows(3)
\end{lstlisting}

\unsection{Joining and stacking tables}{from-pandas:joining}

\textbf{The job.} Put two tables together: side by side by a key (a join), or one under the other
(stacking). In the guides: \href{https://pandas.pydata.org/docs/user_guide/merging.html}{Merge, join,
concatenate and compare};
\href{https://docs.pola.rs/user-guide/transformations/joins/}{Joins} and
\href{https://docs.pola.rs/user-guide/transformations/concatenation/}{Concatenation} in Polars.

\textbf{In simplydata.} \code{join\_with} is \code{merge} and \code{join}. It checks the key for repeats
and says how many rows found a partner, the two most common surprises of a join. \code{keep=} says
which rows stay (\code{how=} in pandas). \code{keep\_rows\_found\_in} and \code{remove\_rows\_found\_in}
are the semi and anti joins. \code{add\_rows\_from} is \code{concat}. \code{show\_differences\_from} is
pandas \code{compare}.

\begin{lstlisting}[style=py]
samples = sdata.get_example("fjord_samples")
stations = sdata.get_example("fjord_stations")
stations.clean_column_names()                # "Station" becomes "station"
samples.join_with(stations, using="station", keep="all rows from this table")
samples.show_first_rows(3)
\end{lstlisting}

\begin{lstlisting}[style=py]
marks = sdata.get_example("marks")
more = sdata.get_example("marks")
more.keep_rows_where("class", is_equal_to="A")
marks.add_rows_from(more)                    # pandas pd.concat([marks, more])
marks.count_rows()
marks.show_duplicate_rows()
\end{lstlisting}

\textbf{Left to pandas and Polars, with a route.}
\begin{itemize}
  \item \textbf{Add two tables cell by cell, with 0 for a missing partner} (pandas
    \code{add(other, fill\_value=0)}). This lines up rows by the index, which simplydata does not have.
    Route: \code{join\_with} on the key, \code{fill\_empty\_cells\_in(..., with\_value=0)}, then
    \code{add\_column} with the sum.
  \item \textbf{Stack and keep only the shared columns} (pandas \code{concat(join="inner")}). Route:
    \code{keep\_columns} on both, then \code{add\_rows\_from}.
  \item \textbf{Put tables side by side by row position} (pandas \code{concat(axis=1)}, Polars
    \code{hstack}). Route: \code{add\_row\_number\_column} on both and \code{join\_with} on it; it is
    safer, because it fails loudly when the tables have different lengths.
  \item \textbf{Join on a condition}, such as a time between a start and an end (Polars
    \code{join\_where}). Route: \code{join\_with(..., keep="every combination")}, then
    \code{keep\_rows\_where} with the condition.
\end{itemize}

\begin{lstlisting}[style=py]
first = sdata.table({"fruit": ["apple", "pear"], "count": [3, 5]})
second = sdata.table({"fruit": ["apple", "plum"], "count": [4, 1]})
second.rename_column("count", to="count_other")
first.join_with(second, using="fruit", keep="all rows from both")
first.fill_empty_cells_in("count", with_value=0)
first.fill_empty_cells_in("count_other", with_value=0)
first.add_column("total", "count + count_other")       # pandas a.add(b, fill_value=0)
first.show()
\end{lstlisting}

\unsection{Reshaping: long and wide}{from-pandas:reshaping}

\textbf{The job.} Turn rows into columns (wide) or columns into rows (long), flip a table, split a cell
into rows, or make yes/no columns. In the guides:
\href{https://pandas.pydata.org/docs/user_guide/reshaping.html}{Reshaping and pivot tables};
\href{https://docs.pola.rs/user-guide/transformations/pivot/}{Pivot} and
\href{https://docs.pola.rs/user-guide/transformations/unpivot/}{Unpivot} in Polars.

\textbf{In simplydata.} \code{spread\_into\_columns} is \code{pivot} (with \code{show=} it is
\code{pivot\_table}); \code{gather\_into\_rows} is \code{melt} and \code{unpivot};
\code{swap\_rows\_and\_columns} is the transpose; \code{split\_into\_rows} is \code{explode};
\code{spread\_into\_yes\_no\_columns} is \code{get\_dummies}. Keep only the columns you need before
spreading: every other column says which row is which.

\begin{lstlisting}[style=py]
table = sdata.get_example("marks")
table.keep_columns("name", "subject", "marks")
table.spread_into_columns(names_from="subject", values_from="marks")   # one row for each child
table.show_first_rows(3)
table.gather_into_rows(["Maths", "Science", "English"], names_to="subject", values_to="marks")
table.show_first_rows(3)
\end{lstlisting}

\begin{lstlisting}[style=py]
fruit = sdata.get_example("fruit")
fruit.keep_columns("day", "fruit", "count")
fruit.spread_into_columns(names_from="fruit", values_from="count", show="total")   # pivot_table
fruit.show()
\end{lstlisting}

\unsection{Text}{from-pandas:text}

\textbf{The job.} Clean, search, cut and join text. pandas keeps these steps under \code{.str}
(\href{https://pandas.pydata.org/docs/user_guide/text.html}{Working with text data}); Polars too
(\href{https://docs.pola.rs/user-guide/expressions/strings/}{Strings}).

\textbf{In simplydata.} The text steps are table words that take the column name:
\code{replace\_text\_in} (\code{str.replace}), \code{remove\_extra\_spaces\_in} (\code{str.strip}),
\code{turn\_into\_lowercase} (\code{str.lower}), \code{add\_column\_with\_part\_of\_text}
(\code{str.slice}, \code{str.extract}), \code{split\_column} (\code{str.split}),
\code{add\_zeros\_in\_front\_of} (\code{str.zfill}), and the conditions \code{contains=},
\code{starts\_with=} and \code{matches\_pattern=} of \code{keep\_rows\_where} (\code{str.contains},
\code{str.match}). In a formula, \code{length(name)} is \code{str.len} and
\code{count\_text(name, ";")} is \code{str.count}.

\begin{lstlisting}[style=py]
samples = sdata.get_example("fjord_samples")
samples.replace_text_in("station", "St.", with_value="Station ")
samples.add_column("crew_size", "count_text(crew, ';') + 1")
samples.keep_rows_where("crew", contains="KL")
samples.add_column_with_part_of_text("first_crew", from_column="crew", before=";")
samples.show_first_rows(3)
\end{lstlisting}

\textbf{A yes/no column from a text match} (pandas \code{np.where(df["name"].str.contains("an"), ...)},
Polars \code{pl.when(pl.col("name").str.contains("an"))}). Rules of \code{add\_group\_column} take
numbers and dates, so count first, then label:

\begin{lstlisting}[style=py]
table = sdata.get_example("marks")
table.add_column("an_count", "count_text(name, 'an')")
table.add_group_column("has_an", rules={"an_count > 0": "yes"}, otherwise="no")
table.count_each("has_an")
\end{lstlisting}

\textbf{Left to pandas and Polars.} Title Case (\code{str.title}), Unicode forms (\code{str.normalize}),
every match of a pattern (\code{str.findall}) and the position of a piece of text (\code{str.find}). The
Unicode form matters when the same letter was typed in two ways, such as an ``{\aa}'' typed on a Mac and
on Windows: the two look the same but are different text, so a join on them fails. The
\code{fjord\_samples} practice table has this on purpose.

\begin{lstlisting}[style=py]
table = sdata.table({"name": ["asha ali", "ben  BERG"], "code": ["S12-x", "S7-y"]})
df = table.as_pandas()
df["name"] = df["name"].str.title().str.normalize("NFC")
df["numbers"] = df["code"].str.findall(r"\d+").str.join(",")
df["dash_at"] = df["code"].str.find("-") + 1          # count from 1, as simplydata does
table = sdata.table(df)
table.show()
\end{lstlisting}

\unsection{Empty cells}{from-pandas:missing}

\textbf{The job.} Find, count, remove and fill empty cells. In the guides:
\href{https://pandas.pydata.org/docs/user_guide/missing_data.html}{Working with missing data};
\href{https://docs.pola.rs/user-guide/expressions/missing-data/}{Missing data} in Polars. pandas mixes
\code{NaN} and \code{None}; Polars keeps \code{null} (empty) apart from \code{NaN} (not a number).
simplydata calls both an empty cell and says so.

\textbf{In simplydata.} \code{count\_empty\_cells\_in} (\code{isna().sum()}),
\code{show\_rows\_with\_empty\_cells\_in}, \code{remove\_rows\_with\_empty\_cells\_in} (\code{dropna}),
\code{fill\_empty\_cells\_in} (\code{fillna}, \code{ffill}, \code{interpolate}) and
\code{add\_is\_empty\_column} (\code{isna}).

\begin{lstlisting}[style=py]
shop = sdata.get_example("shop")
shop.count_empty_cells_in()
shop.fill_empty_cells_in("customers", with_the="line between neighbours", sorted_by="date",
                         for_each_group_in="shop")
shop.fill_empty_cells_in("sales", with_the="value above")
shop.count_empty_cells_in()
\end{lstlisting}

\textbf{Left to pandas and Polars.} Keep rows with at least some filled cells (pandas
\code{dropna(thresh=)}), sort empty cells first (\code{na\_position="first"}), and infinite values
(\code{np.isinf}). They are rare in research tables:

\begin{lstlisting}[style=py]
import numpy as np

table = sdata.table({"a": [1.0, None, np.inf], "b": [None, None, 2.0], "c": [3.0, None, 4.0]})
df = table.as_pandas()
df = df.replace([np.inf, -np.inf], np.nan)       # infinite values become empty
df = df.dropna(thresh=2)                         # keep rows with at least 2 filled cells
table = sdata.table(df)
table.show()
\end{lstlisting}

\unsection{Categories}{from-pandas:categories}

\textbf{The job.} A text column with a fixed set of values in a set order, such as mild, moderate and
severe. In the guides: \href{https://pandas.pydata.org/docs/user_guide/categorical.html}{Categorical
data}; \href{https://docs.pola.rs/user-guide/expressions/categorical-data-and-enums/}{Categorical data
and enums} in Polars.

\textbf{In simplydata.} \code{turn\_into\_categories("column", in\_order=[...])}. Charts and summaries
then follow that order.

\begin{lstlisting}[style=py]
table = sdata.get_example("marks")
table.turn_into_categories("class", in_order=["C", "B", "A"])
table.find_average_of("marks", for_each_group_in="class")    # C, then B, then A
\end{lstlisting}

\textbf{Left to pandas and Polars.} Adding or removing unused categories and reading the codes behind
them (pandas \code{cat.add\_categories}, \code{remove\_unused\_categories}, \code{cat.codes}). simplydata
keeps the list of categories in step with the column for you. For the codes:

\begin{lstlisting}[style=py]
table = sdata.get_example("marks")
table.turn_into_categories("class", in_order=["C", "B", "A"])
df = table.as_pandas()
df["class_code"] = df["class"].cat.codes        # 0 for C, 1 for B, 2 for A
table = sdata.table(df)
table.show_first_rows(3)
\end{lstlisting}

\unsection{Groups}{from-pandas:groups}

\textbf{The job.} Split the rows into groups, sum up each group, or keep some groups. In the guides:
\href{https://pandas.pydata.org/docs/user_guide/groupby.html}{Group by: split-apply-combine};
\href{https://docs.pola.rs/user-guide/expressions/aggregation/}{Aggregation} in Polars.

\textbf{In simplydata.} There is no separate group object. Each summary word takes
\code{for\_each\_group\_in=}: \code{find\_average\_of} (\code{groupby().mean()}),
\code{find\_summary\_of} (\code{agg} with several measures), \code{count\_rows} (\code{size}),
\code{count\_each} (\code{value\_counts}). \code{keep\_groups\_where} is \code{groupby().filter}. A
formula with \code{for\_each\_group\_in=} is \code{transform}.

\begin{lstlisting}[style=py]
table = sdata.get_example("marks")
table.find_average_of("marks", for_each_group_in=["class", "subject"])
summary = table.find_summary_of(["marks", "hours"], show=["average", "count"], for_each_group_in="class")
summary.show()
table.add_column("share_of_class", "marks / total(marks) * 100", for_each_group_in="class")
table.keep_groups_where("average(marks) >= 65", for_each_group_in="class")
table.show_first_rows(3)
\end{lstlisting}

\textbf{The different values of a column} (pandas \code{unique}, Polars \code{unique} on a column).
Keep the column, remove the repeated rows, and take the column. \code{count\_different\_values\_in}
counts them (\code{nunique}).

\begin{lstlisting}[style=py]
table = sdata.get_example("marks")
table.count_different_values_in("name")
table.keep_columns("class")
table.remove_duplicate_rows()
classes = table.get_column("class")      # ["A", "B", "C"], in the order they first appear
print(classes)
\end{lstlisting}

\textbf{Different measures for different columns, one row for each group} (pandas
\code{agg(average=("marks", "mean"), most=("hours", "max"))}, Polars \code{group\_by().agg([...])}),
such as a ``Table 1'' of a paper. \code{find\_summary\_of} gives the same measures for every column.
For different ones, make one table for each measure and join them:

\begin{lstlisting}[style=py]
marks = sdata.get_example("marks")
averages = marks.find_average_of("marks", for_each_group_in="class")
most_hours = marks.find_biggest_of("hours", for_each_group_in="class")
averages.join_with(most_hours, using="class")
averages.show()
\end{lstlisting}

or hand it over:

\begin{lstlisting}[style=py]
table = sdata.get_example("marks")
df = table.as_pandas()
summary = df.groupby("class", as_index=False).agg(average_marks=("marks", "mean"),
                                                  most_hours=("hours", "max"))
summary_table = sdata.table(summary)
summary_table.show()
\end{lstlisting}

\textbf{Join the text of each group into one cell} (pandas \code{agg(";".join)}, Polars
\code{str.join} inside \code{agg}). It is the mirror of \code{split\_into\_rows}, and it is left to
pandas and Polars:

\begin{lstlisting}[style=py]
import polars as pl

lab = sdata.get_example("fjord_lab")
df = lab.as_polars()
genes = df.group_by("sample_id", maintain_order=True).agg(pl.col("gene").str.join(";"))
genes_table = sdata.table(genes)
genes_table.show_first_rows(3)
\end{lstlisting}

\textbf{Left to pandas and Polars.}
\begin{itemize}
  \item \textbf{Your own function on each group} (pandas \code{groupby().apply(f)}, Polars
    \code{map\_groups}). A function per row is \code{add\_column(..., by\_running=f)}; per group, hand
    the table over, as in the example below.
  \item \textbf{A number for each group} (pandas \code{ngroup}). Route:
    \code{add\_rank\_column(..., from\_column="class", biggest\_first=False, when\_tied="no gaps")}.
\end{itemize}

\begin{lstlisting}[style=py]
table = sdata.get_example("marks")
df = table.as_pandas()

def best_subject(group):
    return group.loc[group["marks"].idxmax(), "subject"]

best = df.groupby("name").apply(best_subject, include_groups=False).reset_index(name="best_subject")
best_table = sdata.table(best)
best_table.show_first_rows(3)
\end{lstlisting}

\unsection{Windows: running, rolling and the rows before}{from-pandas:windows}

\textbf{The job.} Work out a value from the rows around each row: a running total, a moving average,
the value of the row before or after. In the guides:
\href{https://pandas.pydata.org/docs/user_guide/window.html}{Windowing operations};
\href{https://docs.pola.rs/user-guide/expressions/window-functions/}{Window functions} in Polars.

\textbf{In simplydata.} \code{add\_running\_column} (\code{cumsum}, \code{cummax},
\code{expanding}), \code{add\_rolling\_column} (\code{rolling}), \code{add\_smoothed\_column}
(\code{ewm} and more), \code{add\_previous\_value\_column} (\code{shift}) and \code{add\_change\_column}
(\code{diff}, \code{pct\_change}). Each takes \code{for\_each\_group\_in=} and \code{sorted\_by=}.

\begin{lstlisting}[style=py]
shop = sdata.get_example("shop")
shop.add_running_column("sales_so_far", from_column="sales", for_each_group_in="shop", sorted_by="date")
shop.add_rolling_column("sales_week", from_column="sales", window=7, show="average",
                        for_each_group_in="shop", sorted_by="date")
shop.add_previous_value_column("sales_yesterday", from_column="sales", for_each_group_in="shop",
                               sorted_by="date")
shop.show_first_rows(3)
\end{lstlisting}

\textbf{The next value} (pandas \code{shift(-1)}, Polars \code{shift(-1)}, R \code{lead}).
\code{add\_previous\_value\_column} looks back. To look ahead, turn the rows upside down, look back,
and turn them again:

\begin{lstlisting}[style=py]
growth = sdata.get_example("growth")
growth.reverse_rows()
growth.add_previous_value_column("next_population", from_column="population", for_each_group_in="country")
growth.reverse_rows()
growth.show_first_rows(3)
\end{lstlisting}

or hand it over:

\begin{lstlisting}[style=py]
growth = sdata.get_example("growth")
df = growth.as_pandas()
df["next_population"] = df.groupby("country")["population"].shift(-1)
growth = sdata.table(df)
growth.show_first_rows(3)
\end{lstlisting}

\textbf{Left to pandas and Polars.} A rolling window that starts with fewer rows (pandas
\code{min\_periods=}, Polars \code{min\_samples=}): simplydata leaves the first cells empty until the
window is full. Runs of the same value (Polars \code{rle\_id}) are rare outside time series:

\begin{lstlisting}[style=py]
import polars as pl

table = sdata.table({"state": ["dry", "dry", "wet", "wet", "wet", "dry"]})
df = table.as_polars().with_columns(pl.col("state").rle_id().alias("run"))
table = sdata.table(df)
table.show()
\end{lstlisting}

\unsection{Dates and times}{from-pandas:dates}

\textbf{The job.} Read dates, take parts of them, round them, add a month, change the time zone and sum
up by week or month. pandas keeps the date steps under \code{.dt}
(\href{https://pandas.pydata.org/docs/user_guide/timeseries.html}{Time series / date functionality},
\href{https://pandas.pydata.org/docs/user_guide/timedeltas.html}{Time deltas}); Polars too
(\href{https://docs.pola.rs/user-guide/transformations/time-series/parsing/}{Parsing},
\href{https://docs.pola.rs/user-guide/transformations/time-series/rolling/}{Grouping},
\href{https://docs.pola.rs/user-guide/transformations/time-series/timezones/}{Time zones}).

\textbf{In simplydata.} \code{turn\_into\_dates} (\code{to\_datetime}),
\code{add\_column\_with\_part\_of\_date} (\code{dt.year}, \code{dt.month}, \code{dt.day\_name}),
\code{round\_dates\_in} (\code{dt.floor}, \code{dt.truncate}), \code{add\_column(..., plus="1 month")}
(\code{DateOffset}, \code{dt.offset\_by}), \code{change\_time\_zone\_in} (\code{tz\_localize} and
\code{tz\_convert}), \code{add\_missing\_dates} (\code{asfreq}) and \code{every=} on every summary word
(\code{resample}, \code{group\_by\_dynamic}).

\begin{lstlisting}[style=py]
shop = sdata.get_example("shop")
shop.turn_into_dates("date")
shop.add_column_with_part_of_date("weekday", from_column="date", part="weekday name")
shop.add_column("check_on", from_column="date", plus="1 month")
shop.find_total_of("sales", for_each_group_in=["shop", "date"], every="month")
\end{lstlisting}

\begin{lstlisting}[style=py]
logger = sdata.get_example("fjord_logger")
logger.turn_into_dates("time_utc")
logger.change_time_zone_in("time_utc", to="Europe/Oslo", recorded_in="UTC")
logger.round_dates_in("time_utc", down_to="hour")
logger.show_first_rows(3)
\end{lstlisting}

\textbf{A date from year, month and day columns} (pandas \code{to\_datetime} of three columns, Polars
\code{pl.date}). Join the parts with a dash, then turn the text into dates:

\begin{lstlisting}[style=py]
table = sdata.table({"year": [2026, 2026], "month": [5, 6], "day": [12, 3]})
table.add_column("date", by_joining_text=["year", "month", "day"], separator="-")
table.turn_into_dates("date", written_like="2026-5-31")
table.show()
\end{lstlisting}

\textbf{Left to pandas and Polars.}
\begin{itemize}
  \item \textbf{Periods} such as ``2026-03'' as one month (pandas \code{to\_period}). Use
    \code{every="month"} in the summary words, or the part \code{"year and month"}.
  \item \textbf{Times of day on their own} (``08:30''), and rows at a time of day (pandas
    \code{between\_time}).
  \item \textbf{Working days} (NumPy \code{busday\_count}, Polars \code{business\_day\_count}).
  \item \textbf{Unix time}, seconds since 1970 (pandas \code{to\_datetime(unit="s")}).
\end{itemize}

\begin{lstlisting}[style=py]
import numpy as np
import pandas as pd

table = sdata.table({"start": ["2026-05-04", "2026-05-11"], "end": ["2026-05-15", "2026-05-22"],
                     "logged": [1778140800, 1778745600]})
df = table.as_pandas()
df["working_days"] = np.busday_count(df["start"].to_numpy("datetime64[D]"),
                                     df["end"].to_numpy("datetime64[D]"))
df["logged"] = pd.to_datetime(df["logged"], unit="s")
df["month"] = pd.to_datetime(df["start"]).dt.to_period("M").astype(str)
table = sdata.table(df)
table.show()
\end{lstlisting}

\unsection{Your own functions}{from-pandas:functions}

\textbf{The job.} Run a Python function that no word has. In the guides:
\href{https://pandas.pydata.org/docs/user_guide/user_defined_functions.html}{User-Defined Functions};
\href{https://docs.pola.rs/user-guide/expressions/user-defined-python-functions/}{User-defined Python
functions} in Polars.

\textbf{In simplydata.} \code{add\_column(..., by\_running=f)} runs your function on each row (pandas
\code{apply(axis=1)}, Polars \code{map\_elements}). The function gets the row as a dict.

\begin{lstlisting}[style=py]
def grade(row):
    return "top" if row["marks"] >= 80 else "other"

table = sdata.get_example("marks")
table.add_column("grade", by_running=grade)
table.count_each("grade")
\end{lstlisting}

\textbf{Left to pandas and Polars.} \code{pipe(f)}, a function on the whole table inside a chain.
simplydata words are already a chain of lines; a plain function that takes the table does the same:

\begin{lstlisting}[style=py]
def clean(table):
    table.clean_column_names()
    table.remove_rows_that_are_all_empty()
    return table

stations = clean(sdata.get_example("fjord_stations"))
stations.get_column_names()
\end{lstlisting}

\unsection{Charts and styled tables}{from-pandas:charts}

pandas draws charts with \code{df.plot}
(\href{https://pandas.pydata.org/docs/user_guide/visualization.html}{Chart visualization}). In
simplydata, \code{table.start\_chart()} gives a chart and its \code{draw\_} words; the charts part of
this manual has them all. Styled tables (pandas
\href{https://pandas.pydata.org/docs/user_guide/style.html}{Table Visualization}: colours, highlights,
number formats) are left to pandas; it needs the package jinja2. \code{table.save("table.html")} saves a
plain web page table.

\unsection{Speed, memory and copies}{from-pandas:speed}

\textbf{The job.} Work fast on big tables, keep memory low, and know when a table is a copy. In the
guides: \href{https://pandas.pydata.org/docs/user_guide/enhancingperf.html}{Enhancing performance},
\href{https://pandas.pydata.org/docs/user_guide/scale.html}{Scaling to large datasets},
\href{https://pandas.pydata.org/docs/user_guide/copy_on_write.html}{Copy-on-Write}; the
\href{https://docs.pola.rs/user-guide/concepts/lazy-api/}{Lazy API} in Polars.

\textbf{In simplydata.} \code{sdata.table("file", fast=True)} runs the table on Polars; the words and
the answers stay the same. A word changes its own table only: \code{backup = table.make\_copy()} keeps
the old one, and \code{table.undo()} goes back one step.

\textbf{Left to pandas and Polars.} Lazy plans and streaming (Polars \code{LazyFrame},
\code{collect(engine="streaming")}), memory use and smaller types (see
\hyperref[from-pandas:kinds]{Kinds of columns}), and sparse storage. Hand the table over for them:

\begin{lstlisting}[style=py]
import polars as pl

river = sdata.get_example("river")
plan = river.as_polars().lazy().filter(pl.col("nitrate") > 3).group_by("station").agg(pl.len())
result = sdata.table(plan.collect())
result.sort_by("station")
result.show()
\end{lstlisting}

\textbf{The numbers as a NumPy array} (pandas and Polars \code{to\_numpy}). Keep the number columns and
hand them over:

\begin{lstlisting}[style=py]
table = sdata.get_example("marks")
table.keep_columns(of_kind="number")
grid = table.as_pandas().to_numpy()
print(grid.shape)
\end{lstlisting}

\textbf{A correlation matrix} (pandas and Polars \code{corr}). This is a statistics word: name the
columns, and \code{find\_correlation\_between} gives every pair with its test.

\begin{lstlisting}[style=py]
table = sdata.get_example("marks")
table.keep_rows_where("subject", is_equal_to="Maths")      # one row for each child
table.find_correlation_between("hours", "marks", "height")
\end{lstlisting}

\unsection{SQL}{from-pandas:sql}

Polars runs SQL on its tables (\href{https://docs.pola.rs/user-guide/sql/intro/}{SQL}); pandas reads
SQL from a database with \code{read\_sql}. In simplydata, \code{table.run\_sql} runs a query on the
table, which is called \code{data} inside the query, and \code{sdata.run\_sql} runs one on files. DuckDB
does the work (pip install ``simplydata[sql]'').

% needs: duckdb
\begin{lstlisting}[style=py]
marks = sdata.get_example("marks")
best = marks.run_sql("SELECT class, count(*) AS children, max(marks) AS best FROM data "
                     "WHERE subject = 'Maths' GROUP BY class ORDER BY class")
best.show()
\end{lstlisting}

\unsection{The jobs left out, in one list}{from-pandas:left-out}

simplydata leaves these jobs to pandas and Polars on purpose. Each is shown above with the hand-over.
\begin{description}[style=nextline,leftmargin=1.2em]
  \item[The row index] set and reset an index, choose and slice rows by label, a MultiIndex, sort by the
    index, repeated labels, and lining up two tables by the index (\code{add(fill\_value=0)}).
    \hyperref[from-pandas:choosing]{Choosing rows and columns}, \hyperref[from-pandas:joining]{Joining}.
  \item[Functions in a chain] choose rows with a function (\code{loc[lambda d: ...]}), your own function on
    each group (\code{groupby().apply}), and \code{pipe}. \hyperref[from-pandas:groups]{Groups},
    \hyperref[from-pandas:functions]{Your own functions}.
  \item[Storage] memory use, smaller types, codes of categories, periods, lazy plans and streaming.
    \hyperref[from-pandas:kinds]{Kinds of columns}, \hyperref[from-pandas:categories]{Categories},
    \hyperref[from-pandas:dates]{Dates and times}, \hyperref[from-pandas:speed]{Speed}.
  \item[Other packages] the numbers as a NumPy array; the correlation matrix is a statistics word.
    \hyperref[from-pandas:speed]{Speed, memory and copies}.
\end{description}
The Polars guide has a page for readers who move the other way, from pandas to Polars:
\href{https://docs.pola.rs/user-guide/migration/pandas/}{Coming from pandas}.

\paragraph{How complete is simplydata?}
We typed every public name of these parts of the library into \code{sdata.help()} (the versions of November 2026).
pandas: DataFrame, Series, groupby, \code{.str}, \code{.dt}, the functions and classes at the top level, and the
rolling, expanding, ewm and resample objects, 800 names. 670 lead to a simplydata word or option, 93 hand over to
pandas with a note, and 37 are left to pandas on purpose because simplydata has no row index and no \code{pipe}.
Polars: DataFrame, Series, expressions, \code{.str}, \code{.dt}, \code{.list}, group\_by and the functions and
classes at the top level, 958 names. 787 lead to a word or option, 98 hand over to Polars with a note (such as its
exponentially weighted spread, folds, struct and JSON columns, engines and settings), and 73 are left to Polars on
purpose: list columns (simplydata keeps one value in each cell) and lazy plans. Every name has an answer in
\code{sdata.help()}.
).
% The topics follow the order of the pandas User Guide and the Polars User Guide, used as a checklist only: the text
% and the examples are our own, and each section links to the official page. Every example in a py box is run by
% tests/test_from_pandas_chapter.py (with "import simplydata as sdata" in front), so it cannot go out of date.
% A line "% needs: duckdb" just before a box says the box needs that package (pip install "simplydata[sql]").

\unchapter{From pandas (and Polars) to simplydata}{from-pandas}

This chapter is for you if you already know pandas or Polars, or if you want to learn them after
simplydata. It walks through the topics of the two official guides, in their order:
the \href{https://pandas.pydata.org/docs/user_guide/index.html}{pandas User Guide} and the
\href{https://docs.pola.rs/user-guide/getting-started/}{Polars User Guide}. For each topic it says
what the job is, which simplydata words do it, and a short example on the practice tables. Some jobs
are left to pandas and Polars on purpose. For those, the chapter shows how to hand the table over and
how to come back.

Every example starts with \code{import simplydata as sdata}. It is written once here, not in each
box. The examples use the practice tables of \code{sdata.get\_example}, so you can run them as they
are.

\unsection{Five things to know first}{from-pandas:first}

\begin{enumerate}
  \item \textbf{A word changes the table.} In pandas and Polars a method gives back a new table, so you
    write \code{df = df.dropna()}. In simplydata, \code{table.remove\_rows\_with\_empty\_cells\_in()}
    changes \code{table} itself, and \code{table.undo()} takes the step back. When you want to keep the
    old table too, make a copy first with \code{table.make\_copy()}. Words that find an answer
    (\code{find\_average\_of}, \code{find\_summary\_of}, \code{run\_sql} and the other
    \code{find\_} words) leave the table as it is: they say the answer and give back a new table with it,
    as in \code{averages = table.find\_average\_of("marks", for\_each\_group\_in="class")}.
  \item \textbf{Rows are numbered from 1.} pandas \code{iloc} and Polars \code{row} count from 0.
    In simplydata the first row is row 1, as people count.
  \item \textbf{There is no row index.} A table is only its columns. Polars works the same way. Where
    pandas uses the index (to select, to sort, to line up two tables), simplydata uses a column.
  \item \textbf{Each word says what it did.} Most words print a sentence: how many rows were kept,
    which rows did not match in a join, what the average is. Read it; it often answers the next
    question.
  \item \textbf{The door is always open.} \code{table.as\_pandas()} gives a pandas DataFrame and
    \code{table.as\_polars()} a Polars DataFrame. \code{sdata.table(df)} turns either one back into a
    simplydata table. Use this for any job this chapter leaves to pandas or Polars.
\end{enumerate}

\begin{lstlisting}[style=py]
table = sdata.get_example("marks")
df = table.as_pandas()                   # a pandas DataFrame: use any pandas code here
df["marks_out_of_10"] = df["marks"] / 10
table = sdata.table(df)                  # back to simplydata
table.show_first_rows(3)
\end{lstlisting}

Do you know a pandas or Polars name, but not the simplydata word? Ask: \code{sdata.help("groupby")}.
Names inside the text and date steps work too: \code{sdata.help("str.replace")},
\code{sdata.help("dt.year")} and \code{sdata.help("groupby.filter")} each name the simplydata word.
The full list is in the section ``Coming from another tool?'' of the simplydata dictionary
(\href{https://github.com/bioforkgithub/simplydata/blob/main/docs/manual/simplydataDictionary.pdf}%
{simplydataDictionary.pdf}).

\unsection{A first look at a table}{from-pandas:first-look}

\textbf{The job.} Open a table, look at the first rows, see its size, the kinds of its columns and a
short summary. In the guides: \href{https://pandas.pydata.org/docs/user_guide/10min.html}{10 minutes
to pandas}, \href{https://pandas.pydata.org/docs/user_guide/dsintro.html}{Intro to data structures},
\href{https://pandas.pydata.org/docs/user_guide/basics.html}{Essential basic functionality};
\href{https://docs.pola.rs/user-guide/getting-started/}{Getting started} and
\href{https://docs.pola.rs/user-guide/concepts/data-types-and-structures/}{Data types and structures}
in Polars.

\textbf{In simplydata.} \code{show\_first\_rows} is pandas \code{head}; \code{count\_rows} and
\code{count\_columns} are \code{shape}; \code{describe\_columns} is a mix of \code{info} and
\code{describe}; \code{find\_summary\_of} gives the measures you choose.

\begin{lstlisting}[style=py]
table = sdata.get_example("marks")
table.show_first_rows(5)
table.count_rows()
table.count_columns()
table.get_column_names()
table.describe_columns()
table.find_summary_of("marks", show=["average", "median", "standard deviation"])
\end{lstlisting}

A pandas Series is one column. In simplydata, \code{table.get\_column("marks")} gives the column as a
list of numbers with its own words (\code{find\_average}, \code{find\_percentile} and more), or as a
plain Python list when the column holds text.

\begin{lstlisting}[style=py]
table = sdata.get_example("marks")
marks = table.get_column("marks")
marks.find_average()
names = table.get_column("name")         # text: a plain Python list
print(names[:3])
\end{lstlisting}

\unsection{Kinds of columns}{from-pandas:kinds}

\textbf{The job.} Make a column hold numbers, text, dates or categories. pandas calls this
\code{astype} and \code{to\_numeric}; Polars calls it \code{cast}
(\href{https://docs.pola.rs/user-guide/expressions/casting/}{Casting}).

\textbf{In simplydata.} One word for each kind: \code{turn\_into\_numbers},
\code{turn\_into\_text}, \code{turn\_into\_dates} and \code{turn\_into\_categories}. Each says how
many cells could not be turned and became empty, so nothing is lost without a word.

\begin{lstlisting}[style=py]
table = sdata.get_example("shop")
table.turn_into_dates("date")
table.turn_into_text("customers")
table.turn_into_numbers("customers")
table.describe_columns()
\end{lstlisting}

\textbf{Left to pandas and Polars.} simplydata chooses the storage type for you (64-bit numbers,
pandas' nullable types, Arrow text). Two jobs about storage are left out on purpose:
\begin{itemize}
  \item \textbf{Smaller types to save memory} (pandas \code{to\_numeric(downcast=)}, Polars
    \code{shrink\_dtype}). A research table rarely needs it.
  \item \textbf{Memory use} (pandas \code{memory\_usage(deep=True)}, Polars \code{estimated\_size}).
\end{itemize}
Hand the table over when you need them:

\begin{lstlisting}[style=py]
import polars as pl

table = sdata.get_example("river")
df = table.as_polars()
print(df.estimated_size("kb"), "kB")
small = pl.DataFrame([df[name].shrink_dtype() if df[name].dtype.is_numeric() else df[name]
                      for name in df.columns])
print(small.estimated_size("kb"), "kB")
table = sdata.table(small)
\end{lstlisting}

\unsection{Reading and saving files}{from-pandas:files}

\textbf{The job.} Read a file into a table and save a table to a file. pandas has one function for each
file type (\href{https://pandas.pydata.org/docs/user_guide/io.html}{IO tools}); so has Polars
(\href{https://docs.pola.rs/user-guide/io/}{IO}).

\textbf{In simplydata.} One word each way. \code{sdata.table("file")} reads .csv, .tsv, .xlsx,
.parquet, .json, .jsonl, .feather, SPSS (.sav), Stata (.dta), database files and more: the end of the
name says the type. \code{table.save("file")} writes them. \code{sdata.run\_sql} and
\code{table.run\_sql} run SQL with DuckDB.

\begin{lstlisting}[style=py]
table = sdata.get_example("marks")
table.save("marks.csv")
table.save("marks.xlsx", in_sheet="2026")
table.save("marks.parquet")
again = sdata.table("marks.parquet")
again.count_rows()
\end{lstlisting}

% needs: duckdb
\begin{lstlisting}[style=py]
river = sdata.get_example("river")
averages = river.run_sql("SELECT station, avg(nitrate) AS nitrate FROM data GROUP BY station ORDER BY station")
averages.show()
\end{lstlisting}

\textbf{A table for a report (Markdown or LaTeX).} simplydata saves .html tables, but not Markdown
(.md) or LaTeX (.tex). pandas has \code{to\_markdown}, which needs the package tabulate, and
\code{to\_latex}, which needs jinja2. You can make both without extra packages: Polars prints a
Markdown table, and a few lines of Python write a LaTeX one.

\begin{lstlisting}[style=py]
import polars as pl

table = sdata.get_example("marks")
averages = table.find_average_of("marks", for_each_group_in="class")    # one row for each class
df = averages.as_polars()
with pl.Config(tbl_formatting="MARKDOWN", tbl_hide_column_data_types=True,
               tbl_hide_dataframe_shape=True, tbl_rows=-1):
    markdown = str(df)
with open("averages.md", "w", encoding="utf-8") as file:
    file.write(markdown + "\n")
print(markdown)
\end{lstlisting}

\begin{lstlisting}[style=py]
table = sdata.get_example("marks")
averages = table.find_average_of("marks", for_each_group_in="class")
columns = averages.get_column_names()
lines = [r"\begin{tabular}{" + "l" * len(columns) + "}", " & ".join(columns) + r" \\ \hline"]
for row in averages.as_list():            # one dict for each row
    lines.append(" & ".join(str(row[name]) for name in columns) + r" \\")
lines.append(r"\end{tabular}")
with open("averages.tex", "w", encoding="utf-8") as file:
    file.write("\n".join(lines) + "\n")
\end{lstlisting}

\textbf{Left to pandas and Polars.} These reading jobs are rare in research, so the manual covers them
here instead of a word:
\begin{itemize}
  \item \textbf{Skip broken lines} (pandas \code{on\_bad\_lines="skip"}). simplydata stops and names the
    broken line instead, because skipping loses data without a word. If you must skip, read with pandas
    and give the result to \code{sdata.table}.
  \item \textbf{Nested JSON} (pandas \code{json\_normalize}, Polars \code{unnest}).
  \item \textbf{The tables of a web page} (pandas \code{read\_html}). simplydata reads files at a web
    address, but not the tables inside a page.
  \item \textbf{The clipboard} (pandas \code{read\_clipboard}).
  \item \textbf{A file too big for memory}, read in pieces or lazily (pandas \code{chunksize=}, Polars
    \code{scan\_csv} and the \href{https://docs.pola.rs/user-guide/concepts/lazy-api/}{Lazy API}).
    simplydata works on tables that fit in memory; \code{fast=True} makes it use Polars underneath.
\end{itemize}

\begin{lstlisting}[style=py]
import pandas as pd

records = [{"station": "St.1", "depth": {"top": 1, "bottom": 10}},
           {"station": "St.2", "depth": {"top": 2, "bottom": 25}}]
flat = pd.json_normalize(records)        # columns station, depth.top, depth.bottom
table = sdata.table(flat)
table.rename_columns({"depth.top": "top", "depth.bottom": "bottom"})
table.show()
\end{lstlisting}

\unsection{Choosing rows and columns}{from-pandas:choosing}

\textbf{The job.} Keep some rows, some columns, or one cell. In pandas this is
\href{https://pandas.pydata.org/docs/user_guide/indexing.html}{Indexing and selecting data}
(\code{loc}, \code{iloc}, \code{query}); in Polars it is \code{filter} and \code{select}
(\href{https://docs.pola.rs/user-guide/expressions/basic-operations/}{Basic operations} and
\href{https://docs.pola.rs/user-guide/expressions/expression-expansion/}{Expression expansion}).

\textbf{In simplydata.} \code{keep\_rows\_where} keeps rows by a condition (or a written condition in
quotes); \code{keep\_columns} keeps columns by name, by kind or by a part of the name;
\code{keep\_rows\_by\_position} is \code{iloc} with rows from 1; \code{get\_cell\_at} is one cell.

\begin{lstlisting}[style=py]
table = sdata.get_example("marks")
table.keep_rows_where("subject", is_equal_to="Maths")
table.keep_rows_where("marks > 60 and class == 'A'")
table.keep_columns("name", "class", "marks")
table.show()
table.get_cell_at(row=1, column="marks")
\end{lstlisting}

\begin{lstlisting}[style=py]
weather = sdata.get_example("weather")
weather.keep_columns(starts_with="wind")      # pandas df.filter(regex="^wind"), Polars pl.col("^wind.*$")
weather.keep_rows_by_position(1, to=3)        # pandas iloc[0:3]
weather.show()
\end{lstlisting}

\textbf{Left to pandas and Polars: the row index.} simplydata has no row index, so these jobs are left
out on purpose. Each has a route with a column.
\begin{itemize}
  \item \textbf{Set or reset an index, choose rows by label, slice labels} (pandas \code{set\_index},
    \code{reset\_index}, \code{loc["a":"c"]},
    \href{https://pandas.pydata.org/docs/user_guide/advanced.html}{MultiIndex}).
    Route: keep the labels in a column and use \code{keep\_rows\_where}.
  \item \textbf{Sort by the index} (pandas \code{sort\_index}). Route: \code{sort\_by} the column.
  \item \textbf{Repeated labels} (pandas \href{https://pandas.pydata.org/docs/user_guide/duplicates.html}{Duplicate
    Labels}). Route: \code{show\_duplicate\_rows(columns=[...])} on the column.
  \item \textbf{Choose rows with a function inside a chain} (pandas \code{loc[lambda d: ...]}). Each
    simplydata word is its own line, so the table you pass is always the table as it is now; a written
    condition, \code{keep\_rows\_where("marks > hours * 10")}, does the same job.
\end{itemize}
When you need the index itself, hand the table over:

\begin{lstlisting}[style=py]
table = sdata.get_example("weather")
df = table.as_pandas().set_index(["city", "month"])     # a MultiIndex
cold = df.loc["Frostby"].sort_index()                   # rows by label, sorted by the index
table = sdata.table(cold.reset_index())                 # the index becomes a column again
table.show_first_rows(3)
\end{lstlisting}

\unsection{Joining and stacking tables}{from-pandas:joining}

\textbf{The job.} Put two tables together: side by side by a key (a join), or one under the other
(stacking). In the guides: \href{https://pandas.pydata.org/docs/user_guide/merging.html}{Merge, join,
concatenate and compare};
\href{https://docs.pola.rs/user-guide/transformations/joins/}{Joins} and
\href{https://docs.pola.rs/user-guide/transformations/concatenation/}{Concatenation} in Polars.

\textbf{In simplydata.} \code{join\_with} is \code{merge} and \code{join}. It checks the key for repeats
and says how many rows found a partner, the two most common surprises of a join. \code{keep=} says
which rows stay (\code{how=} in pandas). \code{keep\_rows\_found\_in} and \code{remove\_rows\_found\_in}
are the semi and anti joins. \code{add\_rows\_from} is \code{concat}. \code{show\_differences\_from} is
pandas \code{compare}.

\begin{lstlisting}[style=py]
samples = sdata.get_example("fjord_samples")
stations = sdata.get_example("fjord_stations")
stations.clean_column_names()                # "Station" becomes "station"
samples.join_with(stations, using="station", keep="all rows from this table")
samples.show_first_rows(3)
\end{lstlisting}

\begin{lstlisting}[style=py]
marks = sdata.get_example("marks")
more = sdata.get_example("marks")
more.keep_rows_where("class", is_equal_to="A")
marks.add_rows_from(more)                    # pandas pd.concat([marks, more])
marks.count_rows()
marks.show_duplicate_rows()
\end{lstlisting}

\textbf{Left to pandas and Polars, with a route.}
\begin{itemize}
  \item \textbf{Add two tables cell by cell, with 0 for a missing partner} (pandas
    \code{add(other, fill\_value=0)}). This lines up rows by the index, which simplydata does not have.
    Route: \code{join\_with} on the key, \code{fill\_empty\_cells\_in(..., with\_value=0)}, then
    \code{add\_column} with the sum.
  \item \textbf{Stack and keep only the shared columns} (pandas \code{concat(join="inner")}). Route:
    \code{keep\_columns} on both, then \code{add\_rows\_from}.
  \item \textbf{Put tables side by side by row position} (pandas \code{concat(axis=1)}, Polars
    \code{hstack}). Route: \code{add\_row\_number\_column} on both and \code{join\_with} on it; it is
    safer, because it fails loudly when the tables have different lengths.
  \item \textbf{Join on a condition}, such as a time between a start and an end (Polars
    \code{join\_where}). Route: \code{join\_with(..., keep="every combination")}, then
    \code{keep\_rows\_where} with the condition.
\end{itemize}

\begin{lstlisting}[style=py]
first = sdata.table({"fruit": ["apple", "pear"], "count": [3, 5]})
second = sdata.table({"fruit": ["apple", "plum"], "count": [4, 1]})
second.rename_column("count", to="count_other")
first.join_with(second, using="fruit", keep="all rows from both")
first.fill_empty_cells_in("count", with_value=0)
first.fill_empty_cells_in("count_other", with_value=0)
first.add_column("total", "count + count_other")       # pandas a.add(b, fill_value=0)
first.show()
\end{lstlisting}

\unsection{Reshaping: long and wide}{from-pandas:reshaping}

\textbf{The job.} Turn rows into columns (wide) or columns into rows (long), flip a table, split a cell
into rows, or make yes/no columns. In the guides:
\href{https://pandas.pydata.org/docs/user_guide/reshaping.html}{Reshaping and pivot tables};
\href{https://docs.pola.rs/user-guide/transformations/pivot/}{Pivot} and
\href{https://docs.pola.rs/user-guide/transformations/unpivot/}{Unpivot} in Polars.

\textbf{In simplydata.} \code{spread\_into\_columns} is \code{pivot} (with \code{show=} it is
\code{pivot\_table}); \code{gather\_into\_rows} is \code{melt} and \code{unpivot};
\code{swap\_rows\_and\_columns} is the transpose; \code{split\_into\_rows} is \code{explode};
\code{spread\_into\_yes\_no\_columns} is \code{get\_dummies}. Keep only the columns you need before
spreading: every other column says which row is which.

\begin{lstlisting}[style=py]
table = sdata.get_example("marks")
table.keep_columns("name", "subject", "marks")
table.spread_into_columns(names_from="subject", values_from="marks")   # one row for each child
table.show_first_rows(3)
table.gather_into_rows(["Maths", "Science", "English"], names_to="subject", values_to="marks")
table.show_first_rows(3)
\end{lstlisting}

\begin{lstlisting}[style=py]
fruit = sdata.get_example("fruit")
fruit.keep_columns("day", "fruit", "count")
fruit.spread_into_columns(names_from="fruit", values_from="count", show="total")   # pivot_table
fruit.show()
\end{lstlisting}

\unsection{Text}{from-pandas:text}

\textbf{The job.} Clean, search, cut and join text. pandas keeps these steps under \code{.str}
(\href{https://pandas.pydata.org/docs/user_guide/text.html}{Working with text data}); Polars too
(\href{https://docs.pola.rs/user-guide/expressions/strings/}{Strings}).

\textbf{In simplydata.} The text steps are table words that take the column name:
\code{replace\_text\_in} (\code{str.replace}), \code{remove\_extra\_spaces\_in} (\code{str.strip}),
\code{turn\_into\_lowercase} (\code{str.lower}), \code{add\_column\_with\_part\_of\_text}
(\code{str.slice}, \code{str.extract}), \code{split\_column} (\code{str.split}),
\code{add\_zeros\_in\_front\_of} (\code{str.zfill}), and the conditions \code{contains=},
\code{starts\_with=} and \code{matches\_pattern=} of \code{keep\_rows\_where} (\code{str.contains},
\code{str.match}). In a formula, \code{length(name)} is \code{str.len} and
\code{count\_text(name, ";")} is \code{str.count}.

\begin{lstlisting}[style=py]
samples = sdata.get_example("fjord_samples")
samples.replace_text_in("station", "St.", with_value="Station ")
samples.add_column("crew_size", "count_text(crew, ';') + 1")
samples.keep_rows_where("crew", contains="KL")
samples.add_column_with_part_of_text("first_crew", from_column="crew", before=";")
samples.show_first_rows(3)
\end{lstlisting}

\textbf{A yes/no column from a text match} (pandas \code{np.where(df["name"].str.contains("an"), ...)},
Polars \code{pl.when(pl.col("name").str.contains("an"))}). Rules of \code{add\_group\_column} take
numbers and dates, so count first, then label:

\begin{lstlisting}[style=py]
table = sdata.get_example("marks")
table.add_column("an_count", "count_text(name, 'an')")
table.add_group_column("has_an", rules={"an_count > 0": "yes"}, otherwise="no")
table.count_each("has_an")
\end{lstlisting}

\textbf{Left to pandas and Polars.} Title Case (\code{str.title}), Unicode forms (\code{str.normalize}),
every match of a pattern (\code{str.findall}) and the position of a piece of text (\code{str.find}). The
Unicode form matters when the same letter was typed in two ways, such as an ``{\aa}'' typed on a Mac and
on Windows: the two look the same but are different text, so a join on them fails. The
\code{fjord\_samples} practice table has this on purpose.

\begin{lstlisting}[style=py]
table = sdata.table({"name": ["asha ali", "ben  BERG"], "code": ["S12-x", "S7-y"]})
df = table.as_pandas()
df["name"] = df["name"].str.title().str.normalize("NFC")
df["numbers"] = df["code"].str.findall(r"\d+").str.join(",")
df["dash_at"] = df["code"].str.find("-") + 1          # count from 1, as simplydata does
table = sdata.table(df)
table.show()
\end{lstlisting}

\unsection{Empty cells}{from-pandas:missing}

\textbf{The job.} Find, count, remove and fill empty cells. In the guides:
\href{https://pandas.pydata.org/docs/user_guide/missing_data.html}{Working with missing data};
\href{https://docs.pola.rs/user-guide/expressions/missing-data/}{Missing data} in Polars. pandas mixes
\code{NaN} and \code{None}; Polars keeps \code{null} (empty) apart from \code{NaN} (not a number).
simplydata calls both an empty cell and says so.

\textbf{In simplydata.} \code{count\_empty\_cells\_in} (\code{isna().sum()}),
\code{show\_rows\_with\_empty\_cells\_in}, \code{remove\_rows\_with\_empty\_cells\_in} (\code{dropna}),
\code{fill\_empty\_cells\_in} (\code{fillna}, \code{ffill}, \code{interpolate}) and
\code{add\_is\_empty\_column} (\code{isna}).

\begin{lstlisting}[style=py]
shop = sdata.get_example("shop")
shop.count_empty_cells_in()
shop.fill_empty_cells_in("customers", with_the="line between neighbours", sorted_by="date",
                         for_each_group_in="shop")
shop.fill_empty_cells_in("sales", with_the="value above")
shop.count_empty_cells_in()
\end{lstlisting}

\textbf{Left to pandas and Polars.} Keep rows with at least some filled cells (pandas
\code{dropna(thresh=)}), sort empty cells first (\code{na\_position="first"}), and infinite values
(\code{np.isinf}). They are rare in research tables:

\begin{lstlisting}[style=py]
import numpy as np

table = sdata.table({"a": [1.0, None, np.inf], "b": [None, None, 2.0], "c": [3.0, None, 4.0]})
df = table.as_pandas()
df = df.replace([np.inf, -np.inf], np.nan)       # infinite values become empty
df = df.dropna(thresh=2)                         # keep rows with at least 2 filled cells
table = sdata.table(df)
table.show()
\end{lstlisting}

\unsection{Categories}{from-pandas:categories}

\textbf{The job.} A text column with a fixed set of values in a set order, such as mild, moderate and
severe. In the guides: \href{https://pandas.pydata.org/docs/user_guide/categorical.html}{Categorical
data}; \href{https://docs.pola.rs/user-guide/expressions/categorical-data-and-enums/}{Categorical data
and enums} in Polars.

\textbf{In simplydata.} \code{turn\_into\_categories("column", in\_order=[...])}. Charts and summaries
then follow that order.

\begin{lstlisting}[style=py]
table = sdata.get_example("marks")
table.turn_into_categories("class", in_order=["C", "B", "A"])
table.find_average_of("marks", for_each_group_in="class")    # C, then B, then A
\end{lstlisting}

\textbf{Left to pandas and Polars.} Adding or removing unused categories and reading the codes behind
them (pandas \code{cat.add\_categories}, \code{remove\_unused\_categories}, \code{cat.codes}). simplydata
keeps the list of categories in step with the column for you. For the codes:

\begin{lstlisting}[style=py]
table = sdata.get_example("marks")
table.turn_into_categories("class", in_order=["C", "B", "A"])
df = table.as_pandas()
df["class_code"] = df["class"].cat.codes        # 0 for C, 1 for B, 2 for A
table = sdata.table(df)
table.show_first_rows(3)
\end{lstlisting}

\unsection{Groups}{from-pandas:groups}

\textbf{The job.} Split the rows into groups, sum up each group, or keep some groups. In the guides:
\href{https://pandas.pydata.org/docs/user_guide/groupby.html}{Group by: split-apply-combine};
\href{https://docs.pola.rs/user-guide/expressions/aggregation/}{Aggregation} in Polars.

\textbf{In simplydata.} There is no separate group object. Each summary word takes
\code{for\_each\_group\_in=}: \code{find\_average\_of} (\code{groupby().mean()}),
\code{find\_summary\_of} (\code{agg} with several measures), \code{count\_rows} (\code{size}),
\code{count\_each} (\code{value\_counts}). \code{keep\_groups\_where} is \code{groupby().filter}. A
formula with \code{for\_each\_group\_in=} is \code{transform}.

\begin{lstlisting}[style=py]
table = sdata.get_example("marks")
table.find_average_of("marks", for_each_group_in=["class", "subject"])
summary = table.find_summary_of(["marks", "hours"], show=["average", "count"], for_each_group_in="class")
summary.show()
table.add_column("share_of_class", "marks / total(marks) * 100", for_each_group_in="class")
table.keep_groups_where("average(marks) >= 65", for_each_group_in="class")
table.show_first_rows(3)
\end{lstlisting}

\textbf{The different values of a column} (pandas \code{unique}, Polars \code{unique} on a column).
Keep the column, remove the repeated rows, and take the column. \code{count\_different\_values\_in}
counts them (\code{nunique}).

\begin{lstlisting}[style=py]
table = sdata.get_example("marks")
table.count_different_values_in("name")
table.keep_columns("class")
table.remove_duplicate_rows()
classes = table.get_column("class")      # ["A", "B", "C"], in the order they first appear
print(classes)
\end{lstlisting}

\textbf{Different measures for different columns, one row for each group} (pandas
\code{agg(average=("marks", "mean"), most=("hours", "max"))}, Polars \code{group\_by().agg([...])}),
such as a ``Table 1'' of a paper. \code{find\_summary\_of} gives the same measures for every column.
For different ones, make one table for each measure and join them:

\begin{lstlisting}[style=py]
marks = sdata.get_example("marks")
averages = marks.find_average_of("marks", for_each_group_in="class")
most_hours = marks.find_biggest_of("hours", for_each_group_in="class")
averages.join_with(most_hours, using="class")
averages.show()
\end{lstlisting}

or hand it over:

\begin{lstlisting}[style=py]
table = sdata.get_example("marks")
df = table.as_pandas()
summary = df.groupby("class", as_index=False).agg(average_marks=("marks", "mean"),
                                                  most_hours=("hours", "max"))
summary_table = sdata.table(summary)
summary_table.show()
\end{lstlisting}

\textbf{Join the text of each group into one cell} (pandas \code{agg(";".join)}, Polars
\code{str.join} inside \code{agg}). It is the mirror of \code{split\_into\_rows}, and it is left to
pandas and Polars:

\begin{lstlisting}[style=py]
import polars as pl

lab = sdata.get_example("fjord_lab")
df = lab.as_polars()
genes = df.group_by("sample_id", maintain_order=True).agg(pl.col("gene").str.join(";"))
genes_table = sdata.table(genes)
genes_table.show_first_rows(3)
\end{lstlisting}

\textbf{Left to pandas and Polars.}
\begin{itemize}
  \item \textbf{Your own function on each group} (pandas \code{groupby().apply(f)}, Polars
    \code{map\_groups}). A function per row is \code{add\_column(..., by\_running=f)}; per group, hand
    the table over, as in the example below.
  \item \textbf{A number for each group} (pandas \code{ngroup}). Route:
    \code{add\_rank\_column(..., from\_column="class", biggest\_first=False, when\_tied="no gaps")}.
\end{itemize}

\begin{lstlisting}[style=py]
table = sdata.get_example("marks")
df = table.as_pandas()

def best_subject(group):
    return group.loc[group["marks"].idxmax(), "subject"]

best = df.groupby("name").apply(best_subject, include_groups=False).reset_index(name="best_subject")
best_table = sdata.table(best)
best_table.show_first_rows(3)
\end{lstlisting}

\unsection{Windows: running, rolling and the rows before}{from-pandas:windows}

\textbf{The job.} Work out a value from the rows around each row: a running total, a moving average,
the value of the row before or after. In the guides:
\href{https://pandas.pydata.org/docs/user_guide/window.html}{Windowing operations};
\href{https://docs.pola.rs/user-guide/expressions/window-functions/}{Window functions} in Polars.

\textbf{In simplydata.} \code{add\_running\_column} (\code{cumsum}, \code{cummax},
\code{expanding}), \code{add\_rolling\_column} (\code{rolling}), \code{add\_smoothed\_column}
(\code{ewm} and more), \code{add\_previous\_value\_column} (\code{shift}) and \code{add\_change\_column}
(\code{diff}, \code{pct\_change}). Each takes \code{for\_each\_group\_in=} and \code{sorted\_by=}.

\begin{lstlisting}[style=py]
shop = sdata.get_example("shop")
shop.add_running_column("sales_so_far", from_column="sales", for_each_group_in="shop", sorted_by="date")
shop.add_rolling_column("sales_week", from_column="sales", window=7, show="average",
                        for_each_group_in="shop", sorted_by="date")
shop.add_previous_value_column("sales_yesterday", from_column="sales", for_each_group_in="shop",
                               sorted_by="date")
shop.show_first_rows(3)
\end{lstlisting}

\textbf{The next value} (pandas \code{shift(-1)}, Polars \code{shift(-1)}, R \code{lead}).
\code{add\_previous\_value\_column} looks back. To look ahead, turn the rows upside down, look back,
and turn them again:

\begin{lstlisting}[style=py]
growth = sdata.get_example("growth")
growth.reverse_rows()
growth.add_previous_value_column("next_population", from_column="population", for_each_group_in="country")
growth.reverse_rows()
growth.show_first_rows(3)
\end{lstlisting}

or hand it over:

\begin{lstlisting}[style=py]
growth = sdata.get_example("growth")
df = growth.as_pandas()
df["next_population"] = df.groupby("country")["population"].shift(-1)
growth = sdata.table(df)
growth.show_first_rows(3)
\end{lstlisting}

\textbf{Left to pandas and Polars.} A rolling window that starts with fewer rows (pandas
\code{min\_periods=}, Polars \code{min\_samples=}): simplydata leaves the first cells empty until the
window is full. Runs of the same value (Polars \code{rle\_id}) are rare outside time series:

\begin{lstlisting}[style=py]
import polars as pl

table = sdata.table({"state": ["dry", "dry", "wet", "wet", "wet", "dry"]})
df = table.as_polars().with_columns(pl.col("state").rle_id().alias("run"))
table = sdata.table(df)
table.show()
\end{lstlisting}

\unsection{Dates and times}{from-pandas:dates}

\textbf{The job.} Read dates, take parts of them, round them, add a month, change the time zone and sum
up by week or month. pandas keeps the date steps under \code{.dt}
(\href{https://pandas.pydata.org/docs/user_guide/timeseries.html}{Time series / date functionality},
\href{https://pandas.pydata.org/docs/user_guide/timedeltas.html}{Time deltas}); Polars too
(\href{https://docs.pola.rs/user-guide/transformations/time-series/parsing/}{Parsing},
\href{https://docs.pola.rs/user-guide/transformations/time-series/rolling/}{Grouping},
\href{https://docs.pola.rs/user-guide/transformations/time-series/timezones/}{Time zones}).

\textbf{In simplydata.} \code{turn\_into\_dates} (\code{to\_datetime}),
\code{add\_column\_with\_part\_of\_date} (\code{dt.year}, \code{dt.month}, \code{dt.day\_name}),
\code{round\_dates\_in} (\code{dt.floor}, \code{dt.truncate}), \code{add\_column(..., plus="1 month")}
(\code{DateOffset}, \code{dt.offset\_by}), \code{change\_time\_zone\_in} (\code{tz\_localize} and
\code{tz\_convert}), \code{add\_missing\_dates} (\code{asfreq}) and \code{every=} on every summary word
(\code{resample}, \code{group\_by\_dynamic}).

\begin{lstlisting}[style=py]
shop = sdata.get_example("shop")
shop.turn_into_dates("date")
shop.add_column_with_part_of_date("weekday", from_column="date", part="weekday name")
shop.add_column("check_on", from_column="date", plus="1 month")
shop.find_total_of("sales", for_each_group_in=["shop", "date"], every="month")
\end{lstlisting}

\begin{lstlisting}[style=py]
logger = sdata.get_example("fjord_logger")
logger.turn_into_dates("time_utc")
logger.change_time_zone_in("time_utc", to="Europe/Oslo", recorded_in="UTC")
logger.round_dates_in("time_utc", down_to="hour")
logger.show_first_rows(3)
\end{lstlisting}

\textbf{A date from year, month and day columns} (pandas \code{to\_datetime} of three columns, Polars
\code{pl.date}). Join the parts with a dash, then turn the text into dates:

\begin{lstlisting}[style=py]
table = sdata.table({"year": [2026, 2026], "month": [5, 6], "day": [12, 3]})
table.add_column("date", by_joining_text=["year", "month", "day"], separator="-")
table.turn_into_dates("date", written_like="2026-5-31")
table.show()
\end{lstlisting}

\textbf{Left to pandas and Polars.}
\begin{itemize}
  \item \textbf{Periods} such as ``2026-03'' as one month (pandas \code{to\_period}). Use
    \code{every="month"} in the summary words, or the part \code{"year and month"}.
  \item \textbf{Times of day on their own} (``08:30''), and rows at a time of day (pandas
    \code{between\_time}).
  \item \textbf{Working days} (NumPy \code{busday\_count}, Polars \code{business\_day\_count}).
  \item \textbf{Unix time}, seconds since 1970 (pandas \code{to\_datetime(unit="s")}).
\end{itemize}

\begin{lstlisting}[style=py]
import numpy as np
import pandas as pd

table = sdata.table({"start": ["2026-05-04", "2026-05-11"], "end": ["2026-05-15", "2026-05-22"],
                     "logged": [1778140800, 1778745600]})
df = table.as_pandas()
df["working_days"] = np.busday_count(df["start"].to_numpy("datetime64[D]"),
                                     df["end"].to_numpy("datetime64[D]"))
df["logged"] = pd.to_datetime(df["logged"], unit="s")
df["month"] = pd.to_datetime(df["start"]).dt.to_period("M").astype(str)
table = sdata.table(df)
table.show()
\end{lstlisting}

\unsection{Your own functions}{from-pandas:functions}

\textbf{The job.} Run a Python function that no word has. In the guides:
\href{https://pandas.pydata.org/docs/user_guide/user_defined_functions.html}{User-Defined Functions};
\href{https://docs.pola.rs/user-guide/expressions/user-defined-python-functions/}{User-defined Python
functions} in Polars.

\textbf{In simplydata.} \code{add\_column(..., by\_running=f)} runs your function on each row (pandas
\code{apply(axis=1)}, Polars \code{map\_elements}). The function gets the row as a dict.

\begin{lstlisting}[style=py]
def grade(row):
    return "top" if row["marks"] >= 80 else "other"

table = sdata.get_example("marks")
table.add_column("grade", by_running=grade)
table.count_each("grade")
\end{lstlisting}

\textbf{Left to pandas and Polars.} \code{pipe(f)}, a function on the whole table inside a chain.
simplydata words are already a chain of lines; a plain function that takes the table does the same:

\begin{lstlisting}[style=py]
def clean(table):
    table.clean_column_names()
    table.remove_rows_that_are_all_empty()
    return table

stations = clean(sdata.get_example("fjord_stations"))
stations.get_column_names()
\end{lstlisting}

\unsection{Charts and styled tables}{from-pandas:charts}

pandas draws charts with \code{df.plot}
(\href{https://pandas.pydata.org/docs/user_guide/visualization.html}{Chart visualization}). In
simplydata, \code{table.start\_chart()} gives a chart and its \code{draw\_} words; the charts part of
this manual has them all. Styled tables (pandas
\href{https://pandas.pydata.org/docs/user_guide/style.html}{Table Visualization}: colours, highlights,
number formats) are left to pandas; it needs the package jinja2. \code{table.save("table.html")} saves a
plain web page table.

\unsection{Speed, memory and copies}{from-pandas:speed}

\textbf{The job.} Work fast on big tables, keep memory low, and know when a table is a copy. In the
guides: \href{https://pandas.pydata.org/docs/user_guide/enhancingperf.html}{Enhancing performance},
\href{https://pandas.pydata.org/docs/user_guide/scale.html}{Scaling to large datasets},
\href{https://pandas.pydata.org/docs/user_guide/copy_on_write.html}{Copy-on-Write}; the
\href{https://docs.pola.rs/user-guide/concepts/lazy-api/}{Lazy API} in Polars.

\textbf{In simplydata.} \code{sdata.table("file", fast=True)} runs the table on Polars; the words and
the answers stay the same. A word changes its own table only: \code{backup = table.make\_copy()} keeps
the old one, and \code{table.undo()} goes back one step.

\textbf{Left to pandas and Polars.} Lazy plans and streaming (Polars \code{LazyFrame},
\code{collect(engine="streaming")}), memory use and smaller types (see
\hyperref[from-pandas:kinds]{Kinds of columns}), and sparse storage. Hand the table over for them:

\begin{lstlisting}[style=py]
import polars as pl

river = sdata.get_example("river")
plan = river.as_polars().lazy().filter(pl.col("nitrate") > 3).group_by("station").agg(pl.len())
result = sdata.table(plan.collect())
result.sort_by("station")
result.show()
\end{lstlisting}

\textbf{The numbers as a NumPy array} (pandas and Polars \code{to\_numpy}). Keep the number columns and
hand them over:

\begin{lstlisting}[style=py]
table = sdata.get_example("marks")
table.keep_columns(of_kind="number")
grid = table.as_pandas().to_numpy()
print(grid.shape)
\end{lstlisting}

\textbf{A correlation matrix} (pandas and Polars \code{corr}). This is a statistics word: name the
columns, and \code{find\_correlation\_between} gives every pair with its test.

\begin{lstlisting}[style=py]
table = sdata.get_example("marks")
table.keep_rows_where("subject", is_equal_to="Maths")      # one row for each child
table.find_correlation_between("hours", "marks", "height")
\end{lstlisting}

\unsection{SQL}{from-pandas:sql}

Polars runs SQL on its tables (\href{https://docs.pola.rs/user-guide/sql/intro/}{SQL}); pandas reads
SQL from a database with \code{read\_sql}. In simplydata, \code{table.run\_sql} runs a query on the
table, which is called \code{data} inside the query, and \code{sdata.run\_sql} runs one on files. DuckDB
does the work (pip install ``simplydata[sql]'').

% needs: duckdb
\begin{lstlisting}[style=py]
marks = sdata.get_example("marks")
best = marks.run_sql("SELECT class, count(*) AS children, max(marks) AS best FROM data "
                     "WHERE subject = 'Maths' GROUP BY class ORDER BY class")
best.show()
\end{lstlisting}

\unsection{The jobs left out, in one list}{from-pandas:left-out}

simplydata leaves these jobs to pandas and Polars on purpose. Each is shown above with the hand-over.
\begin{description}[style=nextline,leftmargin=1.2em]
  \item[The row index] set and reset an index, choose and slice rows by label, a MultiIndex, sort by the
    index, repeated labels, and lining up two tables by the index (\code{add(fill\_value=0)}).
    \hyperref[from-pandas:choosing]{Choosing rows and columns}, \hyperref[from-pandas:joining]{Joining}.
  \item[Functions in a chain] choose rows with a function (\code{loc[lambda d: ...]}), your own function on
    each group (\code{groupby().apply}), and \code{pipe}. \hyperref[from-pandas:groups]{Groups},
    \hyperref[from-pandas:functions]{Your own functions}.
  \item[Storage] memory use, smaller types, codes of categories, periods, lazy plans and streaming.
    \hyperref[from-pandas:kinds]{Kinds of columns}, \hyperref[from-pandas:categories]{Categories},
    \hyperref[from-pandas:dates]{Dates and times}, \hyperref[from-pandas:speed]{Speed}.
  \item[Other packages] the numbers as a NumPy array; the correlation matrix is a statistics word.
    \hyperref[from-pandas:speed]{Speed, memory and copies}.
\end{description}
The Polars guide has a page for readers who move the other way, from pandas to Polars:
\href{https://docs.pola.rs/user-guide/migration/pandas/}{Coming from pandas}.

\paragraph{How complete is simplydata?}
We typed every public name of these parts of the library into \code{sdata.help()} (the versions of November 2026).
pandas: DataFrame, Series, groupby, \code{.str}, \code{.dt}, the functions and classes at the top level, and the
rolling, expanding, ewm and resample objects, 800 names. 670 lead to a simplydata word or option, 93 hand over to
pandas with a note, and 37 are left to pandas on purpose because simplydata has no row index and no \code{pipe}.
Polars: DataFrame, Series, expressions, \code{.str}, \code{.dt}, \code{.list}, group\_by and the functions and
classes at the top level, 958 names. 787 lead to a word or option, 98 hand over to Polars with a note (such as its
exponentially weighted spread, folds, struct and JSON columns, engines and settings), and 73 are left to Polars on
purpose: list columns (simplydata keeps one value in each cell) and lazy plans. Every name has an answer in
\code{sdata.help()}.
).
% The topics follow the order of the pandas User Guide and the Polars User Guide, used as a checklist only: the text
% and the examples are our own, and each section links to the official page. Every example in a py box is run by
% tests/test_from_pandas_chapter.py (with "import simplydata as sdata" in front), so it cannot go out of date.
% A line "% needs: duckdb" just before a box says the box needs that package (pip install "simplydata[sql]").

\unchapter{From pandas (and Polars) to simplydata}{from-pandas}

This chapter is for you if you already know pandas or Polars, or if you want to learn them after
simplydata. It walks through the topics of the two official guides, in their order:
the \href{https://pandas.pydata.org/docs/user_guide/index.html}{pandas User Guide} and the
\href{https://docs.pola.rs/user-guide/getting-started/}{Polars User Guide}. For each topic it says
what the job is, which simplydata words do it, and a short example on the practice tables. Some jobs
are left to pandas and Polars on purpose. For those, the chapter shows how to hand the table over and
how to come back.

Every example starts with \code{import simplydata as sdata}. It is written once here, not in each
box. The examples use the practice tables of \code{sdata.get\_example}, so you can run them as they
are.

\unsection{Five things to know first}{from-pandas:first}

\begin{enumerate}
  \item \textbf{A word changes the table.} In pandas and Polars a method gives back a new table, so you
    write \code{df = df.dropna()}. In simplydata, \code{table.remove\_rows\_with\_empty\_cells\_in()}
    changes \code{table} itself, and \code{table.undo()} takes the step back. When you want to keep the
    old table too, make a copy first with \code{table.make\_copy()}. Words that find an answer
    (\code{find\_average\_of}, \code{find\_summary\_of}, \code{run\_sql} and the other
    \code{find\_} words) leave the table as it is: they say the answer and give back a new table with it,
    as in \code{averages = table.find\_average\_of("marks", for\_each\_group\_in="class")}.
  \item \textbf{Rows are numbered from 1.} pandas \code{iloc} and Polars \code{row} count from 0.
    In simplydata the first row is row 1, as people count.
  \item \textbf{There is no row index.} A table is only its columns. Polars works the same way. Where
    pandas uses the index (to select, to sort, to line up two tables), simplydata uses a column.
  \item \textbf{Each word says what it did.} Most words print a sentence: how many rows were kept,
    which rows did not match in a join, what the average is. Read it; it often answers the next
    question.
  \item \textbf{The door is always open.} \code{table.as\_pandas()} gives a pandas DataFrame and
    \code{table.as\_polars()} a Polars DataFrame. \code{sdata.table(df)} turns either one back into a
    simplydata table. Use this for any job this chapter leaves to pandas or Polars.
\end{enumerate}

\begin{lstlisting}[style=py]
table = sdata.get_example("marks")
df = table.as_pandas()                   # a pandas DataFrame: use any pandas code here
df["marks_out_of_10"] = df["marks"] / 10
table = sdata.table(df)                  # back to simplydata
table.show_first_rows(3)
\end{lstlisting}

Do you know a pandas or Polars name, but not the simplydata word? Ask: \code{sdata.help("groupby")}.
Names inside the text and date steps work too: \code{sdata.help("str.replace")},
\code{sdata.help("dt.year")} and \code{sdata.help("groupby.filter")} each name the simplydata word.
The full list is in the section ``Coming from another tool?'' of the simplydata dictionary
(\href{https://github.com/bioforkgithub/simplydata/blob/main/docs/manual/simplydataDictionary.pdf}%
{simplydataDictionary.pdf}).

\unsection{A first look at a table}{from-pandas:first-look}

\textbf{The job.} Open a table, look at the first rows, see its size, the kinds of its columns and a
short summary. In the guides: \href{https://pandas.pydata.org/docs/user_guide/10min.html}{10 minutes
to pandas}, \href{https://pandas.pydata.org/docs/user_guide/dsintro.html}{Intro to data structures},
\href{https://pandas.pydata.org/docs/user_guide/basics.html}{Essential basic functionality};
\href{https://docs.pola.rs/user-guide/getting-started/}{Getting started} and
\href{https://docs.pola.rs/user-guide/concepts/data-types-and-structures/}{Data types and structures}
in Polars.

\textbf{In simplydata.} \code{show\_first\_rows} is pandas \code{head}; \code{count\_rows} and
\code{count\_columns} are \code{shape}; \code{describe\_columns} is a mix of \code{info} and
\code{describe}; \code{find\_summary\_of} gives the measures you choose.

\begin{lstlisting}[style=py]
table = sdata.get_example("marks")
table.show_first_rows(5)
table.count_rows()
table.count_columns()
table.get_column_names()
table.describe_columns()
table.find_summary_of("marks", show=["average", "median", "standard deviation"])
\end{lstlisting}

A pandas Series is one column. In simplydata, \code{table.get\_column("marks")} gives the column as a
list of numbers with its own words (\code{find\_average}, \code{find\_percentile} and more), or as a
plain Python list when the column holds text.

\begin{lstlisting}[style=py]
table = sdata.get_example("marks")
marks = table.get_column("marks")
marks.find_average()
names = table.get_column("name")         # text: a plain Python list
print(names[:3])
\end{lstlisting}

\unsection{Kinds of columns}{from-pandas:kinds}

\textbf{The job.} Make a column hold numbers, text, dates or categories. pandas calls this
\code{astype} and \code{to\_numeric}; Polars calls it \code{cast}
(\href{https://docs.pola.rs/user-guide/expressions/casting/}{Casting}).

\textbf{In simplydata.} One word for each kind: \code{turn\_into\_numbers},
\code{turn\_into\_text}, \code{turn\_into\_dates} and \code{turn\_into\_categories}. Each says how
many cells could not be turned and became empty, so nothing is lost without a word.

\begin{lstlisting}[style=py]
table = sdata.get_example("shop")
table.turn_into_dates("date")
table.turn_into_text("customers")
table.turn_into_numbers("customers")
table.describe_columns()
\end{lstlisting}

\textbf{Left to pandas and Polars.} simplydata chooses the storage type for you (64-bit numbers,
pandas' nullable types, Arrow text). Two jobs about storage are left out on purpose:
\begin{itemize}
  \item \textbf{Smaller types to save memory} (pandas \code{to\_numeric(downcast=)}, Polars
    \code{shrink\_dtype}). A research table rarely needs it.
  \item \textbf{Memory use} (pandas \code{memory\_usage(deep=True)}, Polars \code{estimated\_size}).
\end{itemize}
Hand the table over when you need them:

\begin{lstlisting}[style=py]
import polars as pl

table = sdata.get_example("river")
df = table.as_polars()
print(df.estimated_size("kb"), "kB")
small = pl.DataFrame([df[name].shrink_dtype() if df[name].dtype.is_numeric() else df[name]
                      for name in df.columns])
print(small.estimated_size("kb"), "kB")
table = sdata.table(small)
\end{lstlisting}

\unsection{Reading and saving files}{from-pandas:files}

\textbf{The job.} Read a file into a table and save a table to a file. pandas has one function for each
file type (\href{https://pandas.pydata.org/docs/user_guide/io.html}{IO tools}); so has Polars
(\href{https://docs.pola.rs/user-guide/io/}{IO}).

\textbf{In simplydata.} One word each way. \code{sdata.table("file")} reads .csv, .tsv, .xlsx,
.parquet, .json, .jsonl, .feather, SPSS (.sav), Stata (.dta), database files and more: the end of the
name says the type. \code{table.save("file")} writes them. \code{sdata.run\_sql} and
\code{table.run\_sql} run SQL with DuckDB.

\begin{lstlisting}[style=py]
table = sdata.get_example("marks")
table.save("marks.csv")
table.save("marks.xlsx", in_sheet="2026")
table.save("marks.parquet")
again = sdata.table("marks.parquet")
again.count_rows()
\end{lstlisting}

% needs: duckdb
\begin{lstlisting}[style=py]
river = sdata.get_example("river")
averages = river.run_sql("SELECT station, avg(nitrate) AS nitrate FROM data GROUP BY station ORDER BY station")
averages.show()
\end{lstlisting}

\textbf{A table for a report (Markdown or LaTeX).} simplydata saves .html tables, but not Markdown
(.md) or LaTeX (.tex). pandas has \code{to\_markdown}, which needs the package tabulate, and
\code{to\_latex}, which needs jinja2. You can make both without extra packages: Polars prints a
Markdown table, and a few lines of Python write a LaTeX one.

\begin{lstlisting}[style=py]
import polars as pl

table = sdata.get_example("marks")
averages = table.find_average_of("marks", for_each_group_in="class")    # one row for each class
df = averages.as_polars()
with pl.Config(tbl_formatting="MARKDOWN", tbl_hide_column_data_types=True,
               tbl_hide_dataframe_shape=True, tbl_rows=-1):
    markdown = str(df)
with open("averages.md", "w", encoding="utf-8") as file:
    file.write(markdown + "\n")
print(markdown)
\end{lstlisting}

\begin{lstlisting}[style=py]
table = sdata.get_example("marks")
averages = table.find_average_of("marks", for_each_group_in="class")
columns = averages.get_column_names()
lines = [r"\begin{tabular}{" + "l" * len(columns) + "}", " & ".join(columns) + r" \\ \hline"]
for row in averages.as_list():            # one dict for each row
    lines.append(" & ".join(str(row[name]) for name in columns) + r" \\")
lines.append(r"\end{tabular}")
with open("averages.tex", "w", encoding="utf-8") as file:
    file.write("\n".join(lines) + "\n")
\end{lstlisting}

\textbf{Left to pandas and Polars.} These reading jobs are rare in research, so the manual covers them
here instead of a word:
\begin{itemize}
  \item \textbf{Skip broken lines} (pandas \code{on\_bad\_lines="skip"}). simplydata stops and names the
    broken line instead, because skipping loses data without a word. If you must skip, read with pandas
    and give the result to \code{sdata.table}.
  \item \textbf{Nested JSON} (pandas \code{json\_normalize}, Polars \code{unnest}).
  \item \textbf{The tables of a web page} (pandas \code{read\_html}). simplydata reads files at a web
    address, but not the tables inside a page.
  \item \textbf{The clipboard} (pandas \code{read\_clipboard}).
  \item \textbf{A file too big for memory}, read in pieces or lazily (pandas \code{chunksize=}, Polars
    \code{scan\_csv} and the \href{https://docs.pola.rs/user-guide/concepts/lazy-api/}{Lazy API}).
    simplydata works on tables that fit in memory; \code{fast=True} makes it use Polars underneath.
\end{itemize}

\begin{lstlisting}[style=py]
import pandas as pd

records = [{"station": "St.1", "depth": {"top": 1, "bottom": 10}},
           {"station": "St.2", "depth": {"top": 2, "bottom": 25}}]
flat = pd.json_normalize(records)        # columns station, depth.top, depth.bottom
table = sdata.table(flat)
table.rename_columns({"depth.top": "top", "depth.bottom": "bottom"})
table.show()
\end{lstlisting}

\unsection{Choosing rows and columns}{from-pandas:choosing}

\textbf{The job.} Keep some rows, some columns, or one cell. In pandas this is
\href{https://pandas.pydata.org/docs/user_guide/indexing.html}{Indexing and selecting data}
(\code{loc}, \code{iloc}, \code{query}); in Polars it is \code{filter} and \code{select}
(\href{https://docs.pola.rs/user-guide/expressions/basic-operations/}{Basic operations} and
\href{https://docs.pola.rs/user-guide/expressions/expression-expansion/}{Expression expansion}).

\textbf{In simplydata.} \code{keep\_rows\_where} keeps rows by a condition (or a written condition in
quotes); \code{keep\_columns} keeps columns by name, by kind or by a part of the name;
\code{keep\_rows\_by\_position} is \code{iloc} with rows from 1; \code{get\_cell\_at} is one cell.

\begin{lstlisting}[style=py]
table = sdata.get_example("marks")
table.keep_rows_where("subject", is_equal_to="Maths")
table.keep_rows_where("marks > 60 and class == 'A'")
table.keep_columns("name", "class", "marks")
table.show()
table.get_cell_at(row=1, column="marks")
\end{lstlisting}

\begin{lstlisting}[style=py]
weather = sdata.get_example("weather")
weather.keep_columns(starts_with="wind")      # pandas df.filter(regex="^wind"), Polars pl.col("^wind.*$")
weather.keep_rows_by_position(1, to=3)        # pandas iloc[0:3]
weather.show()
\end{lstlisting}

\textbf{Left to pandas and Polars: the row index.} simplydata has no row index, so these jobs are left
out on purpose. Each has a route with a column.
\begin{itemize}
  \item \textbf{Set or reset an index, choose rows by label, slice labels} (pandas \code{set\_index},
    \code{reset\_index}, \code{loc["a":"c"]},
    \href{https://pandas.pydata.org/docs/user_guide/advanced.html}{MultiIndex}).
    Route: keep the labels in a column and use \code{keep\_rows\_where}.
  \item \textbf{Sort by the index} (pandas \code{sort\_index}). Route: \code{sort\_by} the column.
  \item \textbf{Repeated labels} (pandas \href{https://pandas.pydata.org/docs/user_guide/duplicates.html}{Duplicate
    Labels}). Route: \code{show\_duplicate\_rows(columns=[...])} on the column.
  \item \textbf{Choose rows with a function inside a chain} (pandas \code{loc[lambda d: ...]}). Each
    simplydata word is its own line, so the table you pass is always the table as it is now; a written
    condition, \code{keep\_rows\_where("marks > hours * 10")}, does the same job.
\end{itemize}
When you need the index itself, hand the table over:

\begin{lstlisting}[style=py]
table = sdata.get_example("weather")
df = table.as_pandas().set_index(["city", "month"])     # a MultiIndex
cold = df.loc["Frostby"].sort_index()                   # rows by label, sorted by the index
table = sdata.table(cold.reset_index())                 # the index becomes a column again
table.show_first_rows(3)
\end{lstlisting}

\unsection{Joining and stacking tables}{from-pandas:joining}

\textbf{The job.} Put two tables together: side by side by a key (a join), or one under the other
(stacking). In the guides: \href{https://pandas.pydata.org/docs/user_guide/merging.html}{Merge, join,
concatenate and compare};
\href{https://docs.pola.rs/user-guide/transformations/joins/}{Joins} and
\href{https://docs.pola.rs/user-guide/transformations/concatenation/}{Concatenation} in Polars.

\textbf{In simplydata.} \code{join\_with} is \code{merge} and \code{join}. It checks the key for repeats
and says how many rows found a partner, the two most common surprises of a join. \code{keep=} says
which rows stay (\code{how=} in pandas). \code{keep\_rows\_found\_in} and \code{remove\_rows\_found\_in}
are the semi and anti joins. \code{add\_rows\_from} is \code{concat}. \code{show\_differences\_from} is
pandas \code{compare}.

\begin{lstlisting}[style=py]
samples = sdata.get_example("fjord_samples")
stations = sdata.get_example("fjord_stations")
stations.clean_column_names()                # "Station" becomes "station"
samples.join_with(stations, using="station", keep="all rows from this table")
samples.show_first_rows(3)
\end{lstlisting}

\begin{lstlisting}[style=py]
marks = sdata.get_example("marks")
more = sdata.get_example("marks")
more.keep_rows_where("class", is_equal_to="A")
marks.add_rows_from(more)                    # pandas pd.concat([marks, more])
marks.count_rows()
marks.show_duplicate_rows()
\end{lstlisting}

\textbf{Left to pandas and Polars, with a route.}
\begin{itemize}
  \item \textbf{Add two tables cell by cell, with 0 for a missing partner} (pandas
    \code{add(other, fill\_value=0)}). This lines up rows by the index, which simplydata does not have.
    Route: \code{join\_with} on the key, \code{fill\_empty\_cells\_in(..., with\_value=0)}, then
    \code{add\_column} with the sum.
  \item \textbf{Stack and keep only the shared columns} (pandas \code{concat(join="inner")}). Route:
    \code{keep\_columns} on both, then \code{add\_rows\_from}.
  \item \textbf{Put tables side by side by row position} (pandas \code{concat(axis=1)}, Polars
    \code{hstack}). Route: \code{add\_row\_number\_column} on both and \code{join\_with} on it; it is
    safer, because it fails loudly when the tables have different lengths.
  \item \textbf{Join on a condition}, such as a time between a start and an end (Polars
    \code{join\_where}). Route: \code{join\_with(..., keep="every combination")}, then
    \code{keep\_rows\_where} with the condition.
\end{itemize}

\begin{lstlisting}[style=py]
first = sdata.table({"fruit": ["apple", "pear"], "count": [3, 5]})
second = sdata.table({"fruit": ["apple", "plum"], "count": [4, 1]})
second.rename_column("count", to="count_other")
first.join_with(second, using="fruit", keep="all rows from both")
first.fill_empty_cells_in("count", with_value=0)
first.fill_empty_cells_in("count_other", with_value=0)
first.add_column("total", "count + count_other")       # pandas a.add(b, fill_value=0)
first.show()
\end{lstlisting}

\unsection{Reshaping: long and wide}{from-pandas:reshaping}

\textbf{The job.} Turn rows into columns (wide) or columns into rows (long), flip a table, split a cell
into rows, or make yes/no columns. In the guides:
\href{https://pandas.pydata.org/docs/user_guide/reshaping.html}{Reshaping and pivot tables};
\href{https://docs.pola.rs/user-guide/transformations/pivot/}{Pivot} and
\href{https://docs.pola.rs/user-guide/transformations/unpivot/}{Unpivot} in Polars.

\textbf{In simplydata.} \code{spread\_into\_columns} is \code{pivot} (with \code{show=} it is
\code{pivot\_table}); \code{gather\_into\_rows} is \code{melt} and \code{unpivot};
\code{swap\_rows\_and\_columns} is the transpose; \code{split\_into\_rows} is \code{explode};
\code{spread\_into\_yes\_no\_columns} is \code{get\_dummies}. Keep only the columns you need before
spreading: every other column says which row is which.

\begin{lstlisting}[style=py]
table = sdata.get_example("marks")
table.keep_columns("name", "subject", "marks")
table.spread_into_columns(names_from="subject", values_from="marks")   # one row for each child
table.show_first_rows(3)
table.gather_into_rows(["Maths", "Science", "English"], names_to="subject", values_to="marks")
table.show_first_rows(3)
\end{lstlisting}

\begin{lstlisting}[style=py]
fruit = sdata.get_example("fruit")
fruit.keep_columns("day", "fruit", "count")
fruit.spread_into_columns(names_from="fruit", values_from="count", show="total")   # pivot_table
fruit.show()
\end{lstlisting}

\unsection{Text}{from-pandas:text}

\textbf{The job.} Clean, search, cut and join text. pandas keeps these steps under \code{.str}
(\href{https://pandas.pydata.org/docs/user_guide/text.html}{Working with text data}); Polars too
(\href{https://docs.pola.rs/user-guide/expressions/strings/}{Strings}).

\textbf{In simplydata.} The text steps are table words that take the column name:
\code{replace\_text\_in} (\code{str.replace}), \code{remove\_extra\_spaces\_in} (\code{str.strip}),
\code{turn\_into\_lowercase} (\code{str.lower}), \code{add\_column\_with\_part\_of\_text}
(\code{str.slice}, \code{str.extract}), \code{split\_column} (\code{str.split}),
\code{add\_zeros\_in\_front\_of} (\code{str.zfill}), and the conditions \code{contains=},
\code{starts\_with=} and \code{matches\_pattern=} of \code{keep\_rows\_where} (\code{str.contains},
\code{str.match}). In a formula, \code{length(name)} is \code{str.len} and
\code{count\_text(name, ";")} is \code{str.count}.

\begin{lstlisting}[style=py]
samples = sdata.get_example("fjord_samples")
samples.replace_text_in("station", "St.", with_value="Station ")
samples.add_column("crew_size", "count_text(crew, ';') + 1")
samples.keep_rows_where("crew", contains="KL")
samples.add_column_with_part_of_text("first_crew", from_column="crew", before=";")
samples.show_first_rows(3)
\end{lstlisting}

\textbf{A yes/no column from a text match} (pandas \code{np.where(df["name"].str.contains("an"), ...)},
Polars \code{pl.when(pl.col("name").str.contains("an"))}). Rules of \code{add\_group\_column} take
numbers and dates, so count first, then label:

\begin{lstlisting}[style=py]
table = sdata.get_example("marks")
table.add_column("an_count", "count_text(name, 'an')")
table.add_group_column("has_an", rules={"an_count > 0": "yes"}, otherwise="no")
table.count_each("has_an")
\end{lstlisting}

\textbf{Left to pandas and Polars.} Title Case (\code{str.title}), Unicode forms (\code{str.normalize}),
every match of a pattern (\code{str.findall}) and the position of a piece of text (\code{str.find}). The
Unicode form matters when the same letter was typed in two ways, such as an ``{\aa}'' typed on a Mac and
on Windows: the two look the same but are different text, so a join on them fails. The
\code{fjord\_samples} practice table has this on purpose.

\begin{lstlisting}[style=py]
table = sdata.table({"name": ["asha ali", "ben  BERG"], "code": ["S12-x", "S7-y"]})
df = table.as_pandas()
df["name"] = df["name"].str.title().str.normalize("NFC")
df["numbers"] = df["code"].str.findall(r"\d+").str.join(",")
df["dash_at"] = df["code"].str.find("-") + 1          # count from 1, as simplydata does
table = sdata.table(df)
table.show()
\end{lstlisting}

\unsection{Empty cells}{from-pandas:missing}

\textbf{The job.} Find, count, remove and fill empty cells. In the guides:
\href{https://pandas.pydata.org/docs/user_guide/missing_data.html}{Working with missing data};
\href{https://docs.pola.rs/user-guide/expressions/missing-data/}{Missing data} in Polars. pandas mixes
\code{NaN} and \code{None}; Polars keeps \code{null} (empty) apart from \code{NaN} (not a number).
simplydata calls both an empty cell and says so.

\textbf{In simplydata.} \code{count\_empty\_cells\_in} (\code{isna().sum()}),
\code{show\_rows\_with\_empty\_cells\_in}, \code{remove\_rows\_with\_empty\_cells\_in} (\code{dropna}),
\code{fill\_empty\_cells\_in} (\code{fillna}, \code{ffill}, \code{interpolate}) and
\code{add\_is\_empty\_column} (\code{isna}).

\begin{lstlisting}[style=py]
shop = sdata.get_example("shop")
shop.count_empty_cells_in()
shop.fill_empty_cells_in("customers", with_the="line between neighbours", sorted_by="date",
                         for_each_group_in="shop")
shop.fill_empty_cells_in("sales", with_the="value above")
shop.count_empty_cells_in()
\end{lstlisting}

\textbf{Left to pandas and Polars.} Keep rows with at least some filled cells (pandas
\code{dropna(thresh=)}), sort empty cells first (\code{na\_position="first"}), and infinite values
(\code{np.isinf}). They are rare in research tables:

\begin{lstlisting}[style=py]
import numpy as np

table = sdata.table({"a": [1.0, None, np.inf], "b": [None, None, 2.0], "c": [3.0, None, 4.0]})
df = table.as_pandas()
df = df.replace([np.inf, -np.inf], np.nan)       # infinite values become empty
df = df.dropna(thresh=2)                         # keep rows with at least 2 filled cells
table = sdata.table(df)
table.show()
\end{lstlisting}

\unsection{Categories}{from-pandas:categories}

\textbf{The job.} A text column with a fixed set of values in a set order, such as mild, moderate and
severe. In the guides: \href{https://pandas.pydata.org/docs/user_guide/categorical.html}{Categorical
data}; \href{https://docs.pola.rs/user-guide/expressions/categorical-data-and-enums/}{Categorical data
and enums} in Polars.

\textbf{In simplydata.} \code{turn\_into\_categories("column", in\_order=[...])}. Charts and summaries
then follow that order.

\begin{lstlisting}[style=py]
table = sdata.get_example("marks")
table.turn_into_categories("class", in_order=["C", "B", "A"])
table.find_average_of("marks", for_each_group_in="class")    # C, then B, then A
\end{lstlisting}

\textbf{Left to pandas and Polars.} Adding or removing unused categories and reading the codes behind
them (pandas \code{cat.add\_categories}, \code{remove\_unused\_categories}, \code{cat.codes}). simplydata
keeps the list of categories in step with the column for you. For the codes:

\begin{lstlisting}[style=py]
table = sdata.get_example("marks")
table.turn_into_categories("class", in_order=["C", "B", "A"])
df = table.as_pandas()
df["class_code"] = df["class"].cat.codes        # 0 for C, 1 for B, 2 for A
table = sdata.table(df)
table.show_first_rows(3)
\end{lstlisting}

\unsection{Groups}{from-pandas:groups}

\textbf{The job.} Split the rows into groups, sum up each group, or keep some groups. In the guides:
\href{https://pandas.pydata.org/docs/user_guide/groupby.html}{Group by: split-apply-combine};
\href{https://docs.pola.rs/user-guide/expressions/aggregation/}{Aggregation} in Polars.

\textbf{In simplydata.} There is no separate group object. Each summary word takes
\code{for\_each\_group\_in=}: \code{find\_average\_of} (\code{groupby().mean()}),
\code{find\_summary\_of} (\code{agg} with several measures), \code{count\_rows} (\code{size}),
\code{count\_each} (\code{value\_counts}). \code{keep\_groups\_where} is \code{groupby().filter}. A
formula with \code{for\_each\_group\_in=} is \code{transform}.

\begin{lstlisting}[style=py]
table = sdata.get_example("marks")
table.find_average_of("marks", for_each_group_in=["class", "subject"])
summary = table.find_summary_of(["marks", "hours"], show=["average", "count"], for_each_group_in="class")
summary.show()
table.add_column("share_of_class", "marks / total(marks) * 100", for_each_group_in="class")
table.keep_groups_where("average(marks) >= 65", for_each_group_in="class")
table.show_first_rows(3)
\end{lstlisting}

\textbf{The different values of a column} (pandas \code{unique}, Polars \code{unique} on a column).
Keep the column, remove the repeated rows, and take the column. \code{count\_different\_values\_in}
counts them (\code{nunique}).

\begin{lstlisting}[style=py]
table = sdata.get_example("marks")
table.count_different_values_in("name")
table.keep_columns("class")
table.remove_duplicate_rows()
classes = table.get_column("class")      # ["A", "B", "C"], in the order they first appear
print(classes)
\end{lstlisting}

\textbf{Different measures for different columns, one row for each group} (pandas
\code{agg(average=("marks", "mean"), most=("hours", "max"))}, Polars \code{group\_by().agg([...])}),
such as a ``Table 1'' of a paper. \code{find\_summary\_of} gives the same measures for every column.
For different ones, make one table for each measure and join them:

\begin{lstlisting}[style=py]
marks = sdata.get_example("marks")
averages = marks.find_average_of("marks", for_each_group_in="class")
most_hours = marks.find_biggest_of("hours", for_each_group_in="class")
averages.join_with(most_hours, using="class")
averages.show()
\end{lstlisting}

or hand it over:

\begin{lstlisting}[style=py]
table = sdata.get_example("marks")
df = table.as_pandas()
summary = df.groupby("class", as_index=False).agg(average_marks=("marks", "mean"),
                                                  most_hours=("hours", "max"))
summary_table = sdata.table(summary)
summary_table.show()
\end{lstlisting}

\textbf{Join the text of each group into one cell} (pandas \code{agg(";".join)}, Polars
\code{str.join} inside \code{agg}). It is the mirror of \code{split\_into\_rows}, and it is left to
pandas and Polars:

\begin{lstlisting}[style=py]
import polars as pl

lab = sdata.get_example("fjord_lab")
df = lab.as_polars()
genes = df.group_by("sample_id", maintain_order=True).agg(pl.col("gene").str.join(";"))
genes_table = sdata.table(genes)
genes_table.show_first_rows(3)
\end{lstlisting}

\textbf{Left to pandas and Polars.}
\begin{itemize}
  \item \textbf{Your own function on each group} (pandas \code{groupby().apply(f)}, Polars
    \code{map\_groups}). A function per row is \code{add\_column(..., by\_running=f)}; per group, hand
    the table over, as in the example below.
  \item \textbf{A number for each group} (pandas \code{ngroup}). Route:
    \code{add\_rank\_column(..., from\_column="class", biggest\_first=False, when\_tied="no gaps")}.
\end{itemize}

\begin{lstlisting}[style=py]
table = sdata.get_example("marks")
df = table.as_pandas()

def best_subject(group):
    return group.loc[group["marks"].idxmax(), "subject"]

best = df.groupby("name").apply(best_subject, include_groups=False).reset_index(name="best_subject")
best_table = sdata.table(best)
best_table.show_first_rows(3)
\end{lstlisting}

\unsection{Windows: running, rolling and the rows before}{from-pandas:windows}

\textbf{The job.} Work out a value from the rows around each row: a running total, a moving average,
the value of the row before or after. In the guides:
\href{https://pandas.pydata.org/docs/user_guide/window.html}{Windowing operations};
\href{https://docs.pola.rs/user-guide/expressions/window-functions/}{Window functions} in Polars.

\textbf{In simplydata.} \code{add\_running\_column} (\code{cumsum}, \code{cummax},
\code{expanding}), \code{add\_rolling\_column} (\code{rolling}), \code{add\_smoothed\_column}
(\code{ewm} and more), \code{add\_previous\_value\_column} (\code{shift}) and \code{add\_change\_column}
(\code{diff}, \code{pct\_change}). Each takes \code{for\_each\_group\_in=} and \code{sorted\_by=}.

\begin{lstlisting}[style=py]
shop = sdata.get_example("shop")
shop.add_running_column("sales_so_far", from_column="sales", for_each_group_in="shop", sorted_by="date")
shop.add_rolling_column("sales_week", from_column="sales", window=7, show="average",
                        for_each_group_in="shop", sorted_by="date")
shop.add_previous_value_column("sales_yesterday", from_column="sales", for_each_group_in="shop",
                               sorted_by="date")
shop.show_first_rows(3)
\end{lstlisting}

\textbf{The next value} (pandas \code{shift(-1)}, Polars \code{shift(-1)}, R \code{lead}).
\code{add\_previous\_value\_column} looks back. To look ahead, turn the rows upside down, look back,
and turn them again:

\begin{lstlisting}[style=py]
growth = sdata.get_example("growth")
growth.reverse_rows()
growth.add_previous_value_column("next_population", from_column="population", for_each_group_in="country")
growth.reverse_rows()
growth.show_first_rows(3)
\end{lstlisting}

or hand it over:

\begin{lstlisting}[style=py]
growth = sdata.get_example("growth")
df = growth.as_pandas()
df["next_population"] = df.groupby("country")["population"].shift(-1)
growth = sdata.table(df)
growth.show_first_rows(3)
\end{lstlisting}

\textbf{Left to pandas and Polars.} A rolling window that starts with fewer rows (pandas
\code{min\_periods=}, Polars \code{min\_samples=}): simplydata leaves the first cells empty until the
window is full. Runs of the same value (Polars \code{rle\_id}) are rare outside time series:

\begin{lstlisting}[style=py]
import polars as pl

table = sdata.table({"state": ["dry", "dry", "wet", "wet", "wet", "dry"]})
df = table.as_polars().with_columns(pl.col("state").rle_id().alias("run"))
table = sdata.table(df)
table.show()
\end{lstlisting}

\unsection{Dates and times}{from-pandas:dates}

\textbf{The job.} Read dates, take parts of them, round them, add a month, change the time zone and sum
up by week or month. pandas keeps the date steps under \code{.dt}
(\href{https://pandas.pydata.org/docs/user_guide/timeseries.html}{Time series / date functionality},
\href{https://pandas.pydata.org/docs/user_guide/timedeltas.html}{Time deltas}); Polars too
(\href{https://docs.pola.rs/user-guide/transformations/time-series/parsing/}{Parsing},
\href{https://docs.pola.rs/user-guide/transformations/time-series/rolling/}{Grouping},
\href{https://docs.pola.rs/user-guide/transformations/time-series/timezones/}{Time zones}).

\textbf{In simplydata.} \code{turn\_into\_dates} (\code{to\_datetime}),
\code{add\_column\_with\_part\_of\_date} (\code{dt.year}, \code{dt.month}, \code{dt.day\_name}),
\code{round\_dates\_in} (\code{dt.floor}, \code{dt.truncate}), \code{add\_column(..., plus="1 month")}
(\code{DateOffset}, \code{dt.offset\_by}), \code{change\_time\_zone\_in} (\code{tz\_localize} and
\code{tz\_convert}), \code{add\_missing\_dates} (\code{asfreq}) and \code{every=} on every summary word
(\code{resample}, \code{group\_by\_dynamic}).

\begin{lstlisting}[style=py]
shop = sdata.get_example("shop")
shop.turn_into_dates("date")
shop.add_column_with_part_of_date("weekday", from_column="date", part="weekday name")
shop.add_column("check_on", from_column="date", plus="1 month")
shop.find_total_of("sales", for_each_group_in=["shop", "date"], every="month")
\end{lstlisting}

\begin{lstlisting}[style=py]
logger = sdata.get_example("fjord_logger")
logger.turn_into_dates("time_utc")
logger.change_time_zone_in("time_utc", to="Europe/Oslo", recorded_in="UTC")
logger.round_dates_in("time_utc", down_to="hour")
logger.show_first_rows(3)
\end{lstlisting}

\textbf{A date from year, month and day columns} (pandas \code{to\_datetime} of three columns, Polars
\code{pl.date}). Join the parts with a dash, then turn the text into dates:

\begin{lstlisting}[style=py]
table = sdata.table({"year": [2026, 2026], "month": [5, 6], "day": [12, 3]})
table.add_column("date", by_joining_text=["year", "month", "day"], separator="-")
table.turn_into_dates("date", written_like="2026-5-31")
table.show()
\end{lstlisting}

\textbf{Left to pandas and Polars.}
\begin{itemize}
  \item \textbf{Periods} such as ``2026-03'' as one month (pandas \code{to\_period}). Use
    \code{every="month"} in the summary words, or the part \code{"year and month"}.
  \item \textbf{Times of day on their own} (``08:30''), and rows at a time of day (pandas
    \code{between\_time}).
  \item \textbf{Working days} (NumPy \code{busday\_count}, Polars \code{business\_day\_count}).
  \item \textbf{Unix time}, seconds since 1970 (pandas \code{to\_datetime(unit="s")}).
\end{itemize}

\begin{lstlisting}[style=py]
import numpy as np
import pandas as pd

table = sdata.table({"start": ["2026-05-04", "2026-05-11"], "end": ["2026-05-15", "2026-05-22"],
                     "logged": [1778140800, 1778745600]})
df = table.as_pandas()
df["working_days"] = np.busday_count(df["start"].to_numpy("datetime64[D]"),
                                     df["end"].to_numpy("datetime64[D]"))
df["logged"] = pd.to_datetime(df["logged"], unit="s")
df["month"] = pd.to_datetime(df["start"]).dt.to_period("M").astype(str)
table = sdata.table(df)
table.show()
\end{lstlisting}

\unsection{Your own functions}{from-pandas:functions}

\textbf{The job.} Run a Python function that no word has. In the guides:
\href{https://pandas.pydata.org/docs/user_guide/user_defined_functions.html}{User-Defined Functions};
\href{https://docs.pola.rs/user-guide/expressions/user-defined-python-functions/}{User-defined Python
functions} in Polars.

\textbf{In simplydata.} \code{add\_column(..., by\_running=f)} runs your function on each row (pandas
\code{apply(axis=1)}, Polars \code{map\_elements}). The function gets the row as a dict.

\begin{lstlisting}[style=py]
def grade(row):
    return "top" if row["marks"] >= 80 else "other"

table = sdata.get_example("marks")
table.add_column("grade", by_running=grade)
table.count_each("grade")
\end{lstlisting}

\textbf{Left to pandas and Polars.} \code{pipe(f)}, a function on the whole table inside a chain.
simplydata words are already a chain of lines; a plain function that takes the table does the same:

\begin{lstlisting}[style=py]
def clean(table):
    table.clean_column_names()
    table.remove_rows_that_are_all_empty()
    return table

stations = clean(sdata.get_example("fjord_stations"))
stations.get_column_names()
\end{lstlisting}

\unsection{Charts and styled tables}{from-pandas:charts}

pandas draws charts with \code{df.plot}
(\href{https://pandas.pydata.org/docs/user_guide/visualization.html}{Chart visualization}). In
simplydata, \code{table.start\_chart()} gives a chart and its \code{draw\_} words; the charts part of
this manual has them all. Styled tables (pandas
\href{https://pandas.pydata.org/docs/user_guide/style.html}{Table Visualization}: colours, highlights,
number formats) are left to pandas; it needs the package jinja2. \code{table.save("table.html")} saves a
plain web page table.

\unsection{Speed, memory and copies}{from-pandas:speed}

\textbf{The job.} Work fast on big tables, keep memory low, and know when a table is a copy. In the
guides: \href{https://pandas.pydata.org/docs/user_guide/enhancingperf.html}{Enhancing performance},
\href{https://pandas.pydata.org/docs/user_guide/scale.html}{Scaling to large datasets},
\href{https://pandas.pydata.org/docs/user_guide/copy_on_write.html}{Copy-on-Write}; the
\href{https://docs.pola.rs/user-guide/concepts/lazy-api/}{Lazy API} in Polars.

\textbf{In simplydata.} \code{sdata.table("file", fast=True)} runs the table on Polars; the words and
the answers stay the same. A word changes its own table only: \code{backup = table.make\_copy()} keeps
the old one, and \code{table.undo()} goes back one step.

\textbf{Left to pandas and Polars.} Lazy plans and streaming (Polars \code{LazyFrame},
\code{collect(engine="streaming")}), memory use and smaller types (see
\hyperref[from-pandas:kinds]{Kinds of columns}), and sparse storage. Hand the table over for them:

\begin{lstlisting}[style=py]
import polars as pl

river = sdata.get_example("river")
plan = river.as_polars().lazy().filter(pl.col("nitrate") > 3).group_by("station").agg(pl.len())
result = sdata.table(plan.collect())
result.sort_by("station")
result.show()
\end{lstlisting}

\textbf{The numbers as a NumPy array} (pandas and Polars \code{to\_numpy}). Keep the number columns and
hand them over:

\begin{lstlisting}[style=py]
table = sdata.get_example("marks")
table.keep_columns(of_kind="number")
grid = table.as_pandas().to_numpy()
print(grid.shape)
\end{lstlisting}

\textbf{A correlation matrix} (pandas and Polars \code{corr}). This is a statistics word: name the
columns, and \code{find\_correlation\_between} gives every pair with its test.

\begin{lstlisting}[style=py]
table = sdata.get_example("marks")
table.keep_rows_where("subject", is_equal_to="Maths")      # one row for each child
table.find_correlation_between("hours", "marks", "height")
\end{lstlisting}

\unsection{SQL}{from-pandas:sql}

Polars runs SQL on its tables (\href{https://docs.pola.rs/user-guide/sql/intro/}{SQL}); pandas reads
SQL from a database with \code{read\_sql}. In simplydata, \code{table.run\_sql} runs a query on the
table, which is called \code{data} inside the query, and \code{sdata.run\_sql} runs one on files. DuckDB
does the work (pip install ``simplydata[sql]'').

% needs: duckdb
\begin{lstlisting}[style=py]
marks = sdata.get_example("marks")
best = marks.run_sql("SELECT class, count(*) AS children, max(marks) AS best FROM data "
                     "WHERE subject = 'Maths' GROUP BY class ORDER BY class")
best.show()
\end{lstlisting}

\unsection{The jobs left out, in one list}{from-pandas:left-out}

simplydata leaves these jobs to pandas and Polars on purpose. Each is shown above with the hand-over.
\begin{description}[style=nextline,leftmargin=1.2em]
  \item[The row index] set and reset an index, choose and slice rows by label, a MultiIndex, sort by the
    index, repeated labels, and lining up two tables by the index (\code{add(fill\_value=0)}).
    \hyperref[from-pandas:choosing]{Choosing rows and columns}, \hyperref[from-pandas:joining]{Joining}.
  \item[Functions in a chain] choose rows with a function (\code{loc[lambda d: ...]}), your own function on
    each group (\code{groupby().apply}), and \code{pipe}. \hyperref[from-pandas:groups]{Groups},
    \hyperref[from-pandas:functions]{Your own functions}.
  \item[Storage] memory use, smaller types, codes of categories, periods, lazy plans and streaming.
    \hyperref[from-pandas:kinds]{Kinds of columns}, \hyperref[from-pandas:categories]{Categories},
    \hyperref[from-pandas:dates]{Dates and times}, \hyperref[from-pandas:speed]{Speed}.
  \item[Other packages] the numbers as a NumPy array; the correlation matrix is a statistics word.
    \hyperref[from-pandas:speed]{Speed, memory and copies}.
\end{description}
The Polars guide has a page for readers who move the other way, from pandas to Polars:
\href{https://docs.pola.rs/user-guide/migration/pandas/}{Coming from pandas}.

\paragraph{How complete is simplydata?}
We typed every public name of these parts of the library into \code{sdata.help()} (the versions of November 2026).
pandas: DataFrame, Series, groupby, \code{.str}, \code{.dt}, the functions and classes at the top level, and the
rolling, expanding, ewm and resample objects, 800 names. 670 lead to a simplydata word or option, 93 hand over to
pandas with a note, and 37 are left to pandas on purpose because simplydata has no row index and no \code{pipe}.
Polars: DataFrame, Series, expressions, \code{.str}, \code{.dt}, \code{.list}, group\_by and the functions and
classes at the top level, 958 names. 787 lead to a word or option, 98 hand over to Polars with a note (such as its
exponentially weighted spread, folds, struct and JSON columns, engines and settings), and 73 are left to Polars on
purpose: list columns (simplydata keeps one value in each cell) and lazy plans. Every name has an answer in
\code{sdata.help()}.
).
% The topics follow the order of the pandas User Guide and the Polars User Guide, used as a checklist only: the text
% and the examples are our own, and each section links to the official page. Every example in a py box is run by
% tests/test_from_pandas_chapter.py (with "import simplydata as sdata" in front), so it cannot go out of date.
% A line "% needs: duckdb" just before a box says the box needs that package (pip install "simplydata[sql]").

\unchapter{From pandas (and Polars) to simplydata}{from-pandas}

This chapter is for you if you already know pandas or Polars, or if you want to learn them after
simplydata. It walks through the topics of the two official guides, in their order:
the \href{https://pandas.pydata.org/docs/user_guide/index.html}{pandas User Guide} and the
\href{https://docs.pola.rs/user-guide/getting-started/}{Polars User Guide}. For each topic it says
what the job is, which simplydata words do it, and a short example on the practice tables. Some jobs
are left to pandas and Polars on purpose. For those, the chapter shows how to hand the table over and
how to come back.

Every example starts with \code{import simplydata as sdata}. It is written once here, not in each
box. The examples use the practice tables of \code{sdata.get\_example}, so you can run them as they
are.

\unsection{Five things to know first}{from-pandas:first}

\begin{enumerate}
  \item \textbf{A word changes the table.} In pandas and Polars a method gives back a new table, so you
    write \code{df = df.dropna()}. In simplydata, \code{table.remove\_rows\_with\_empty\_cells\_in()}
    changes \code{table} itself, and \code{table.undo()} takes the step back. When you want to keep the
    old table too, make a copy first with \code{table.make\_copy()}. Words that find an answer
    (\code{find\_average\_of}, \code{find\_summary\_of}, \code{run\_sql} and the other
    \code{find\_} words) leave the table as it is: they say the answer and give back a new table with it,
    as in \code{averages = table.find\_average\_of("marks", for\_each\_group\_in="class")}.
  \item \textbf{Rows are numbered from 1.} pandas \code{iloc} and Polars \code{row} count from 0.
    In simplydata the first row is row 1, as people count.
  \item \textbf{There is no row index.} A table is only its columns. Polars works the same way. Where
    pandas uses the index (to select, to sort, to line up two tables), simplydata uses a column.
  \item \textbf{Each word says what it did.} Most words print a sentence: how many rows were kept,
    which rows did not match in a join, what the average is. Read it; it often answers the next
    question.
  \item \textbf{The door is always open.} \code{table.as\_pandas()} gives a pandas DataFrame and
    \code{table.as\_polars()} a Polars DataFrame. \code{sdata.table(df)} turns either one back into a
    simplydata table. Use this for any job this chapter leaves to pandas or Polars.
\end{enumerate}

\begin{lstlisting}[style=py]
table = sdata.get_example("marks")
df = table.as_pandas()                   # a pandas DataFrame: use any pandas code here
df["marks_out_of_10"] = df["marks"] / 10
table = sdata.table(df)                  # back to simplydata
table.show_first_rows(3)
\end{lstlisting}

Do you know a pandas or Polars name, but not the simplydata word? Ask: \code{sdata.help("groupby")}.
Names inside the text and date steps work too: \code{sdata.help("str.replace")},
\code{sdata.help("dt.year")} and \code{sdata.help("groupby.filter")} each name the simplydata word.
The full list is in the section ``Coming from another tool?'' of the simplydata dictionary
(\href{https://github.com/bioforkgithub/simplydata/blob/main/docs/manual/simplydataDictionary.pdf}%
{simplydataDictionary.pdf}).

\unsection{A first look at a table}{from-pandas:first-look}

\textbf{The job.} Open a table, look at the first rows, see its size, the kinds of its columns and a
short summary. In the guides: \href{https://pandas.pydata.org/docs/user_guide/10min.html}{10 minutes
to pandas}, \href{https://pandas.pydata.org/docs/user_guide/dsintro.html}{Intro to data structures},
\href{https://pandas.pydata.org/docs/user_guide/basics.html}{Essential basic functionality};
\href{https://docs.pola.rs/user-guide/getting-started/}{Getting started} and
\href{https://docs.pola.rs/user-guide/concepts/data-types-and-structures/}{Data types and structures}
in Polars.

\textbf{In simplydata.} \code{show\_first\_rows} is pandas \code{head}; \code{count\_rows} and
\code{count\_columns} are \code{shape}; \code{describe\_columns} is a mix of \code{info} and
\code{describe}; \code{find\_summary\_of} gives the measures you choose.

\begin{lstlisting}[style=py]
table = sdata.get_example("marks")
table.show_first_rows(5)
table.count_rows()
table.count_columns()
table.get_column_names()
table.describe_columns()
table.find_summary_of("marks", show=["average", "median", "standard deviation"])
\end{lstlisting}

A pandas Series is one column. In simplydata, \code{table.get\_column("marks")} gives the column as a
list of numbers with its own words (\code{find\_average}, \code{find\_percentile} and more), or as a
plain Python list when the column holds text.

\begin{lstlisting}[style=py]
table = sdata.get_example("marks")
marks = table.get_column("marks")
marks.find_average()
names = table.get_column("name")         # text: a plain Python list
print(names[:3])
\end{lstlisting}

\unsection{Kinds of columns}{from-pandas:kinds}

\textbf{The job.} Make a column hold numbers, text, dates or categories. pandas calls this
\code{astype} and \code{to\_numeric}; Polars calls it \code{cast}
(\href{https://docs.pola.rs/user-guide/expressions/casting/}{Casting}).

\textbf{In simplydata.} One word for each kind: \code{turn\_into\_numbers},
\code{turn\_into\_text}, \code{turn\_into\_dates} and \code{turn\_into\_categories}. Each says how
many cells could not be turned and became empty, so nothing is lost without a word.

\begin{lstlisting}[style=py]
table = sdata.get_example("shop")
table.turn_into_dates("date")
table.turn_into_text("customers")
table.turn_into_numbers("customers")
table.describe_columns()
\end{lstlisting}

\textbf{Left to pandas and Polars.} simplydata chooses the storage type for you (64-bit numbers,
pandas' nullable types, Arrow text). Two jobs about storage are left out on purpose:
\begin{itemize}
  \item \textbf{Smaller types to save memory} (pandas \code{to\_numeric(downcast=)}, Polars
    \code{shrink\_dtype}). A research table rarely needs it.
  \item \textbf{Memory use} (pandas \code{memory\_usage(deep=True)}, Polars \code{estimated\_size}).
\end{itemize}
Hand the table over when you need them:

\begin{lstlisting}[style=py]
import polars as pl

table = sdata.get_example("river")
df = table.as_polars()
print(df.estimated_size("kb"), "kB")
small = pl.DataFrame([df[name].shrink_dtype() if df[name].dtype.is_numeric() else df[name]
                      for name in df.columns])
print(small.estimated_size("kb"), "kB")
table = sdata.table(small)
\end{lstlisting}

\unsection{Reading and saving files}{from-pandas:files}

\textbf{The job.} Read a file into a table and save a table to a file. pandas has one function for each
file type (\href{https://pandas.pydata.org/docs/user_guide/io.html}{IO tools}); so has Polars
(\href{https://docs.pola.rs/user-guide/io/}{IO}).

\textbf{In simplydata.} One word each way. \code{sdata.table("file")} reads .csv, .tsv, .xlsx,
.parquet, .json, .jsonl, .feather, SPSS (.sav), Stata (.dta), database files and more: the end of the
name says the type. \code{table.save("file")} writes them. \code{sdata.run\_sql} and
\code{table.run\_sql} run SQL with DuckDB.

\begin{lstlisting}[style=py]
table = sdata.get_example("marks")
table.save("marks.csv")
table.save("marks.xlsx", in_sheet="2026")
table.save("marks.parquet")
again = sdata.table("marks.parquet")
again.count_rows()
\end{lstlisting}

% needs: duckdb
\begin{lstlisting}[style=py]
river = sdata.get_example("river")
averages = river.run_sql("SELECT station, avg(nitrate) AS nitrate FROM data GROUP BY station ORDER BY station")
averages.show()
\end{lstlisting}

\textbf{A table for a report (Markdown or LaTeX).} simplydata saves .html tables, but not Markdown
(.md) or LaTeX (.tex). pandas has \code{to\_markdown}, which needs the package tabulate, and
\code{to\_latex}, which needs jinja2. You can make both without extra packages: Polars prints a
Markdown table, and a few lines of Python write a LaTeX one.

\begin{lstlisting}[style=py]
import polars as pl

table = sdata.get_example("marks")
averages = table.find_average_of("marks", for_each_group_in="class")    # one row for each class
df = averages.as_polars()
with pl.Config(tbl_formatting="MARKDOWN", tbl_hide_column_data_types=True,
               tbl_hide_dataframe_shape=True, tbl_rows=-1):
    markdown = str(df)
with open("averages.md", "w", encoding="utf-8") as file:
    file.write(markdown + "\n")
print(markdown)
\end{lstlisting}

\begin{lstlisting}[style=py]
table = sdata.get_example("marks")
averages = table.find_average_of("marks", for_each_group_in="class")
columns = averages.get_column_names()
lines = [r"\begin{tabular}{" + "l" * len(columns) + "}", " & ".join(columns) + r" \\ \hline"]
for row in averages.as_list():            # one dict for each row
    lines.append(" & ".join(str(row[name]) for name in columns) + r" \\")
lines.append(r"\end{tabular}")
with open("averages.tex", "w", encoding="utf-8") as file:
    file.write("\n".join(lines) + "\n")
\end{lstlisting}

\textbf{Left to pandas and Polars.} These reading jobs are rare in research, so the manual covers them
here instead of a word:
\begin{itemize}
  \item \textbf{Skip broken lines} (pandas \code{on\_bad\_lines="skip"}). simplydata stops and names the
    broken line instead, because skipping loses data without a word. If you must skip, read with pandas
    and give the result to \code{sdata.table}.
  \item \textbf{Nested JSON} (pandas \code{json\_normalize}, Polars \code{unnest}).
  \item \textbf{The tables of a web page} (pandas \code{read\_html}). simplydata reads files at a web
    address, but not the tables inside a page.
  \item \textbf{The clipboard} (pandas \code{read\_clipboard}).
  \item \textbf{A file too big for memory}, read in pieces or lazily (pandas \code{chunksize=}, Polars
    \code{scan\_csv} and the \href{https://docs.pola.rs/user-guide/concepts/lazy-api/}{Lazy API}).
    simplydata works on tables that fit in memory; \code{fast=True} makes it use Polars underneath.
\end{itemize}

\begin{lstlisting}[style=py]
import pandas as pd

records = [{"station": "St.1", "depth": {"top": 1, "bottom": 10}},
           {"station": "St.2", "depth": {"top": 2, "bottom": 25}}]
flat = pd.json_normalize(records)        # columns station, depth.top, depth.bottom
table = sdata.table(flat)
table.rename_columns({"depth.top": "top", "depth.bottom": "bottom"})
table.show()
\end{lstlisting}

\unsection{Choosing rows and columns}{from-pandas:choosing}

\textbf{The job.} Keep some rows, some columns, or one cell. In pandas this is
\href{https://pandas.pydata.org/docs/user_guide/indexing.html}{Indexing and selecting data}
(\code{loc}, \code{iloc}, \code{query}); in Polars it is \code{filter} and \code{select}
(\href{https://docs.pola.rs/user-guide/expressions/basic-operations/}{Basic operations} and
\href{https://docs.pola.rs/user-guide/expressions/expression-expansion/}{Expression expansion}).

\textbf{In simplydata.} \code{keep\_rows\_where} keeps rows by a condition (or a written condition in
quotes); \code{keep\_columns} keeps columns by name, by kind or by a part of the name;
\code{keep\_rows\_by\_position} is \code{iloc} with rows from 1; \code{get\_cell\_at} is one cell.

\begin{lstlisting}[style=py]
table = sdata.get_example("marks")
table.keep_rows_where("subject", is_equal_to="Maths")
table.keep_rows_where("marks > 60 and class == 'A'")
table.keep_columns("name", "class", "marks")
table.show()
table.get_cell_at(row=1, column="marks")
\end{lstlisting}

\begin{lstlisting}[style=py]
weather = sdata.get_example("weather")
weather.keep_columns(starts_with="wind")      # pandas df.filter(regex="^wind"), Polars pl.col("^wind.*$")
weather.keep_rows_by_position(1, to=3)        # pandas iloc[0:3]
weather.show()
\end{lstlisting}

\textbf{Left to pandas and Polars: the row index.} simplydata has no row index, so these jobs are left
out on purpose. Each has a route with a column.
\begin{itemize}
  \item \textbf{Set or reset an index, choose rows by label, slice labels} (pandas \code{set\_index},
    \code{reset\_index}, \code{loc["a":"c"]},
    \href{https://pandas.pydata.org/docs/user_guide/advanced.html}{MultiIndex}).
    Route: keep the labels in a column and use \code{keep\_rows\_where}.
  \item \textbf{Sort by the index} (pandas \code{sort\_index}). Route: \code{sort\_by} the column.
  \item \textbf{Repeated labels} (pandas \href{https://pandas.pydata.org/docs/user_guide/duplicates.html}{Duplicate
    Labels}). Route: \code{show\_duplicate\_rows(columns=[...])} on the column.
  \item \textbf{Choose rows with a function inside a chain} (pandas \code{loc[lambda d: ...]}). Each
    simplydata word is its own line, so the table you pass is always the table as it is now; a written
    condition, \code{keep\_rows\_where("marks > hours * 10")}, does the same job.
\end{itemize}
When you need the index itself, hand the table over:

\begin{lstlisting}[style=py]
table = sdata.get_example("weather")
df = table.as_pandas().set_index(["city", "month"])     # a MultiIndex
cold = df.loc["Frostby"].sort_index()                   # rows by label, sorted by the index
table = sdata.table(cold.reset_index())                 # the index becomes a column again
table.show_first_rows(3)
\end{lstlisting}

\unsection{Joining and stacking tables}{from-pandas:joining}

\textbf{The job.} Put two tables together: side by side by a key (a join), or one under the other
(stacking). In the guides: \href{https://pandas.pydata.org/docs/user_guide/merging.html}{Merge, join,
concatenate and compare};
\href{https://docs.pola.rs/user-guide/transformations/joins/}{Joins} and
\href{https://docs.pola.rs/user-guide/transformations/concatenation/}{Concatenation} in Polars.

\textbf{In simplydata.} \code{join\_with} is \code{merge} and \code{join}. It checks the key for repeats
and says how many rows found a partner, the two most common surprises of a join. \code{keep=} says
which rows stay (\code{how=} in pandas). \code{keep\_rows\_found\_in} and \code{remove\_rows\_found\_in}
are the semi and anti joins. \code{add\_rows\_from} is \code{concat}. \code{show\_differences\_from} is
pandas \code{compare}.

\begin{lstlisting}[style=py]
samples = sdata.get_example("fjord_samples")
stations = sdata.get_example("fjord_stations")
stations.clean_column_names()                # "Station" becomes "station"
samples.join_with(stations, using="station", keep="all rows from this table")
samples.show_first_rows(3)
\end{lstlisting}

\begin{lstlisting}[style=py]
marks = sdata.get_example("marks")
more = sdata.get_example("marks")
more.keep_rows_where("class", is_equal_to="A")
marks.add_rows_from(more)                    # pandas pd.concat([marks, more])
marks.count_rows()
marks.show_duplicate_rows()
\end{lstlisting}

\textbf{Left to pandas and Polars, with a route.}
\begin{itemize}
  \item \textbf{Add two tables cell by cell, with 0 for a missing partner} (pandas
    \code{add(other, fill\_value=0)}). This lines up rows by the index, which simplydata does not have.
    Route: \code{join\_with} on the key, \code{fill\_empty\_cells\_in(..., with\_value=0)}, then
    \code{add\_column} with the sum.
  \item \textbf{Stack and keep only the shared columns} (pandas \code{concat(join="inner")}). Route:
    \code{keep\_columns} on both, then \code{add\_rows\_from}.
  \item \textbf{Put tables side by side by row position} (pandas \code{concat(axis=1)}, Polars
    \code{hstack}). Route: \code{add\_row\_number\_column} on both and \code{join\_with} on it; it is
    safer, because it fails loudly when the tables have different lengths.
  \item \textbf{Join on a condition}, such as a time between a start and an end (Polars
    \code{join\_where}). Route: \code{join\_with(..., keep="every combination")}, then
    \code{keep\_rows\_where} with the condition.
\end{itemize}

\begin{lstlisting}[style=py]
first = sdata.table({"fruit": ["apple", "pear"], "count": [3, 5]})
second = sdata.table({"fruit": ["apple", "plum"], "count": [4, 1]})
second.rename_column("count", to="count_other")
first.join_with(second, using="fruit", keep="all rows from both")
first.fill_empty_cells_in("count", with_value=0)
first.fill_empty_cells_in("count_other", with_value=0)
first.add_column("total", "count + count_other")       # pandas a.add(b, fill_value=0)
first.show()
\end{lstlisting}

\unsection{Reshaping: long and wide}{from-pandas:reshaping}

\textbf{The job.} Turn rows into columns (wide) or columns into rows (long), flip a table, split a cell
into rows, or make yes/no columns. In the guides:
\href{https://pandas.pydata.org/docs/user_guide/reshaping.html}{Reshaping and pivot tables};
\href{https://docs.pola.rs/user-guide/transformations/pivot/}{Pivot} and
\href{https://docs.pola.rs/user-guide/transformations/unpivot/}{Unpivot} in Polars.

\textbf{In simplydata.} \code{spread\_into\_columns} is \code{pivot} (with \code{show=} it is
\code{pivot\_table}); \code{gather\_into\_rows} is \code{melt} and \code{unpivot};
\code{swap\_rows\_and\_columns} is the transpose; \code{split\_into\_rows} is \code{explode};
\code{spread\_into\_yes\_no\_columns} is \code{get\_dummies}. Keep only the columns you need before
spreading: every other column says which row is which.

\begin{lstlisting}[style=py]
table = sdata.get_example("marks")
table.keep_columns("name", "subject", "marks")
table.spread_into_columns(names_from="subject", values_from="marks")   # one row for each child
table.show_first_rows(3)
table.gather_into_rows(["Maths", "Science", "English"], names_to="subject", values_to="marks")
table.show_first_rows(3)
\end{lstlisting}

\begin{lstlisting}[style=py]
fruit = sdata.get_example("fruit")
fruit.keep_columns("day", "fruit", "count")
fruit.spread_into_columns(names_from="fruit", values_from="count", show="total")   # pivot_table
fruit.show()
\end{lstlisting}

\unsection{Text}{from-pandas:text}

\textbf{The job.} Clean, search, cut and join text. pandas keeps these steps under \code{.str}
(\href{https://pandas.pydata.org/docs/user_guide/text.html}{Working with text data}); Polars too
(\href{https://docs.pola.rs/user-guide/expressions/strings/}{Strings}).

\textbf{In simplydata.} The text steps are table words that take the column name:
\code{replace\_text\_in} (\code{str.replace}), \code{remove\_extra\_spaces\_in} (\code{str.strip}),
\code{turn\_into\_lowercase} (\code{str.lower}), \code{add\_column\_with\_part\_of\_text}
(\code{str.slice}, \code{str.extract}), \code{split\_column} (\code{str.split}),
\code{add\_zeros\_in\_front\_of} (\code{str.zfill}), and the conditions \code{contains=},
\code{starts\_with=} and \code{matches\_pattern=} of \code{keep\_rows\_where} (\code{str.contains},
\code{str.match}). In a formula, \code{length(name)} is \code{str.len} and
\code{count\_text(name, ";")} is \code{str.count}.

\begin{lstlisting}[style=py]
samples = sdata.get_example("fjord_samples")
samples.replace_text_in("station", "St.", with_value="Station ")
samples.add_column("crew_size", "count_text(crew, ';') + 1")
samples.keep_rows_where("crew", contains="KL")
samples.add_column_with_part_of_text("first_crew", from_column="crew", before=";")
samples.show_first_rows(3)
\end{lstlisting}

\textbf{A yes/no column from a text match} (pandas \code{np.where(df["name"].str.contains("an"), ...)},
Polars \code{pl.when(pl.col("name").str.contains("an"))}). Rules of \code{add\_group\_column} take
numbers and dates, so count first, then label:

\begin{lstlisting}[style=py]
table = sdata.get_example("marks")
table.add_column("an_count", "count_text(name, 'an')")
table.add_group_column("has_an", rules={"an_count > 0": "yes"}, otherwise="no")
table.count_each("has_an")
\end{lstlisting}

\textbf{Left to pandas and Polars.} Title Case (\code{str.title}), Unicode forms (\code{str.normalize}),
every match of a pattern (\code{str.findall}) and the position of a piece of text (\code{str.find}). The
Unicode form matters when the same letter was typed in two ways, such as an ``{\aa}'' typed on a Mac and
on Windows: the two look the same but are different text, so a join on them fails. The
\code{fjord\_samples} practice table has this on purpose.

\begin{lstlisting}[style=py]
table = sdata.table({"name": ["asha ali", "ben  BERG"], "code": ["S12-x", "S7-y"]})
df = table.as_pandas()
df["name"] = df["name"].str.title().str.normalize("NFC")
df["numbers"] = df["code"].str.findall(r"\d+").str.join(",")
df["dash_at"] = df["code"].str.find("-") + 1          # count from 1, as simplydata does
table = sdata.table(df)
table.show()
\end{lstlisting}

\unsection{Empty cells}{from-pandas:missing}

\textbf{The job.} Find, count, remove and fill empty cells. In the guides:
\href{https://pandas.pydata.org/docs/user_guide/missing_data.html}{Working with missing data};
\href{https://docs.pola.rs/user-guide/expressions/missing-data/}{Missing data} in Polars. pandas mixes
\code{NaN} and \code{None}; Polars keeps \code{null} (empty) apart from \code{NaN} (not a number).
simplydata calls both an empty cell and says so.

\textbf{In simplydata.} \code{count\_empty\_cells\_in} (\code{isna().sum()}),
\code{show\_rows\_with\_empty\_cells\_in}, \code{remove\_rows\_with\_empty\_cells\_in} (\code{dropna}),
\code{fill\_empty\_cells\_in} (\code{fillna}, \code{ffill}, \code{interpolate}) and
\code{add\_is\_empty\_column} (\code{isna}).

\begin{lstlisting}[style=py]
shop = sdata.get_example("shop")
shop.count_empty_cells_in()
shop.fill_empty_cells_in("customers", with_the="line between neighbours", sorted_by="date",
                         for_each_group_in="shop")
shop.fill_empty_cells_in("sales", with_the="value above")
shop.count_empty_cells_in()
\end{lstlisting}

\textbf{Left to pandas and Polars.} Keep rows with at least some filled cells (pandas
\code{dropna(thresh=)}), sort empty cells first (\code{na\_position="first"}), and infinite values
(\code{np.isinf}). They are rare in research tables:

\begin{lstlisting}[style=py]
import numpy as np

table = sdata.table({"a": [1.0, None, np.inf], "b": [None, None, 2.0], "c": [3.0, None, 4.0]})
df = table.as_pandas()
df = df.replace([np.inf, -np.inf], np.nan)       # infinite values become empty
df = df.dropna(thresh=2)                         # keep rows with at least 2 filled cells
table = sdata.table(df)
table.show()
\end{lstlisting}

\unsection{Categories}{from-pandas:categories}

\textbf{The job.} A text column with a fixed set of values in a set order, such as mild, moderate and
severe. In the guides: \href{https://pandas.pydata.org/docs/user_guide/categorical.html}{Categorical
data}; \href{https://docs.pola.rs/user-guide/expressions/categorical-data-and-enums/}{Categorical data
and enums} in Polars.

\textbf{In simplydata.} \code{turn\_into\_categories("column", in\_order=[...])}. Charts and summaries
then follow that order.

\begin{lstlisting}[style=py]
table = sdata.get_example("marks")
table.turn_into_categories("class", in_order=["C", "B", "A"])
table.find_average_of("marks", for_each_group_in="class")    # C, then B, then A
\end{lstlisting}

\textbf{Left to pandas and Polars.} Adding or removing unused categories and reading the codes behind
them (pandas \code{cat.add\_categories}, \code{remove\_unused\_categories}, \code{cat.codes}). simplydata
keeps the list of categories in step with the column for you. For the codes:

\begin{lstlisting}[style=py]
table = sdata.get_example("marks")
table.turn_into_categories("class", in_order=["C", "B", "A"])
df = table.as_pandas()
df["class_code"] = df["class"].cat.codes        # 0 for C, 1 for B, 2 for A
table = sdata.table(df)
table.show_first_rows(3)
\end{lstlisting}

\unsection{Groups}{from-pandas:groups}

\textbf{The job.} Split the rows into groups, sum up each group, or keep some groups. In the guides:
\href{https://pandas.pydata.org/docs/user_guide/groupby.html}{Group by: split-apply-combine};
\href{https://docs.pola.rs/user-guide/expressions/aggregation/}{Aggregation} in Polars.

\textbf{In simplydata.} There is no separate group object. Each summary word takes
\code{for\_each\_group\_in=}: \code{find\_average\_of} (\code{groupby().mean()}),
\code{find\_summary\_of} (\code{agg} with several measures), \code{count\_rows} (\code{size}),
\code{count\_each} (\code{value\_counts}). \code{keep\_groups\_where} is \code{groupby().filter}. A
formula with \code{for\_each\_group\_in=} is \code{transform}.

\begin{lstlisting}[style=py]
table = sdata.get_example("marks")
table.find_average_of("marks", for_each_group_in=["class", "subject"])
summary = table.find_summary_of(["marks", "hours"], show=["average", "count"], for_each_group_in="class")
summary.show()
table.add_column("share_of_class", "marks / total(marks) * 100", for_each_group_in="class")
table.keep_groups_where("average(marks) >= 65", for_each_group_in="class")
table.show_first_rows(3)
\end{lstlisting}

\textbf{The different values of a column} (pandas \code{unique}, Polars \code{unique} on a column).
Keep the column, remove the repeated rows, and take the column. \code{count\_different\_values\_in}
counts them (\code{nunique}).

\begin{lstlisting}[style=py]
table = sdata.get_example("marks")
table.count_different_values_in("name")
table.keep_columns("class")
table.remove_duplicate_rows()
classes = table.get_column("class")      # ["A", "B", "C"], in the order they first appear
print(classes)
\end{lstlisting}

\textbf{Different measures for different columns, one row for each group} (pandas
\code{agg(average=("marks", "mean"), most=("hours", "max"))}, Polars \code{group\_by().agg([...])}),
such as a ``Table 1'' of a paper. \code{find\_summary\_of} gives the same measures for every column.
For different ones, make one table for each measure and join them:

\begin{lstlisting}[style=py]
marks = sdata.get_example("marks")
averages = marks.find_average_of("marks", for_each_group_in="class")
most_hours = marks.find_biggest_of("hours", for_each_group_in="class")
averages.join_with(most_hours, using="class")
averages.show()
\end{lstlisting}

or hand it over:

\begin{lstlisting}[style=py]
table = sdata.get_example("marks")
df = table.as_pandas()
summary = df.groupby("class", as_index=False).agg(average_marks=("marks", "mean"),
                                                  most_hours=("hours", "max"))
summary_table = sdata.table(summary)
summary_table.show()
\end{lstlisting}

\textbf{Join the text of each group into one cell} (pandas \code{agg(";".join)}, Polars
\code{str.join} inside \code{agg}). It is the mirror of \code{split\_into\_rows}, and it is left to
pandas and Polars:

\begin{lstlisting}[style=py]
import polars as pl

lab = sdata.get_example("fjord_lab")
df = lab.as_polars()
genes = df.group_by("sample_id", maintain_order=True).agg(pl.col("gene").str.join(";"))
genes_table = sdata.table(genes)
genes_table.show_first_rows(3)
\end{lstlisting}

\textbf{Left to pandas and Polars.}
\begin{itemize}
  \item \textbf{Your own function on each group} (pandas \code{groupby().apply(f)}, Polars
    \code{map\_groups}). A function per row is \code{add\_column(..., by\_running=f)}; per group, hand
    the table over, as in the example below.
  \item \textbf{A number for each group} (pandas \code{ngroup}). Route:
    \code{add\_rank\_column(..., from\_column="class", biggest\_first=False, when\_tied="no gaps")}.
\end{itemize}

\begin{lstlisting}[style=py]
table = sdata.get_example("marks")
df = table.as_pandas()

def best_subject(group):
    return group.loc[group["marks"].idxmax(), "subject"]

best = df.groupby("name").apply(best_subject, include_groups=False).reset_index(name="best_subject")
best_table = sdata.table(best)
best_table.show_first_rows(3)
\end{lstlisting}

\unsection{Windows: running, rolling and the rows before}{from-pandas:windows}

\textbf{The job.} Work out a value from the rows around each row: a running total, a moving average,
the value of the row before or after. In the guides:
\href{https://pandas.pydata.org/docs/user_guide/window.html}{Windowing operations};
\href{https://docs.pola.rs/user-guide/expressions/window-functions/}{Window functions} in Polars.

\textbf{In simplydata.} \code{add\_running\_column} (\code{cumsum}, \code{cummax},
\code{expanding}), \code{add\_rolling\_column} (\code{rolling}), \code{add\_smoothed\_column}
(\code{ewm} and more), \code{add\_previous\_value\_column} (\code{shift}) and \code{add\_change\_column}
(\code{diff}, \code{pct\_change}). Each takes \code{for\_each\_group\_in=} and \code{sorted\_by=}.

\begin{lstlisting}[style=py]
shop = sdata.get_example("shop")
shop.add_running_column("sales_so_far", from_column="sales", for_each_group_in="shop", sorted_by="date")
shop.add_rolling_column("sales_week", from_column="sales", window=7, show="average",
                        for_each_group_in="shop", sorted_by="date")
shop.add_previous_value_column("sales_yesterday", from_column="sales", for_each_group_in="shop",
                               sorted_by="date")
shop.show_first_rows(3)
\end{lstlisting}

\textbf{The next value} (pandas \code{shift(-1)}, Polars \code{shift(-1)}, R \code{lead}).
\code{add\_previous\_value\_column} looks back. To look ahead, turn the rows upside down, look back,
and turn them again:

\begin{lstlisting}[style=py]
growth = sdata.get_example("growth")
growth.reverse_rows()
growth.add_previous_value_column("next_population", from_column="population", for_each_group_in="country")
growth.reverse_rows()
growth.show_first_rows(3)
\end{lstlisting}

or hand it over:

\begin{lstlisting}[style=py]
growth = sdata.get_example("growth")
df = growth.as_pandas()
df["next_population"] = df.groupby("country")["population"].shift(-1)
growth = sdata.table(df)
growth.show_first_rows(3)
\end{lstlisting}

\textbf{Left to pandas and Polars.} A rolling window that starts with fewer rows (pandas
\code{min\_periods=}, Polars \code{min\_samples=}): simplydata leaves the first cells empty until the
window is full. Runs of the same value (Polars \code{rle\_id}) are rare outside time series:

\begin{lstlisting}[style=py]
import polars as pl

table = sdata.table({"state": ["dry", "dry", "wet", "wet", "wet", "dry"]})
df = table.as_polars().with_columns(pl.col("state").rle_id().alias("run"))
table = sdata.table(df)
table.show()
\end{lstlisting}

\unsection{Dates and times}{from-pandas:dates}

\textbf{The job.} Read dates, take parts of them, round them, add a month, change the time zone and sum
up by week or month. pandas keeps the date steps under \code{.dt}
(\href{https://pandas.pydata.org/docs/user_guide/timeseries.html}{Time series / date functionality},
\href{https://pandas.pydata.org/docs/user_guide/timedeltas.html}{Time deltas}); Polars too
(\href{https://docs.pola.rs/user-guide/transformations/time-series/parsing/}{Parsing},
\href{https://docs.pola.rs/user-guide/transformations/time-series/rolling/}{Grouping},
\href{https://docs.pola.rs/user-guide/transformations/time-series/timezones/}{Time zones}).

\textbf{In simplydata.} \code{turn\_into\_dates} (\code{to\_datetime}),
\code{add\_column\_with\_part\_of\_date} (\code{dt.year}, \code{dt.month}, \code{dt.day\_name}),
\code{round\_dates\_in} (\code{dt.floor}, \code{dt.truncate}), \code{add\_column(..., plus="1 month")}
(\code{DateOffset}, \code{dt.offset\_by}), \code{change\_time\_zone\_in} (\code{tz\_localize} and
\code{tz\_convert}), \code{add\_missing\_dates} (\code{asfreq}) and \code{every=} on every summary word
(\code{resample}, \code{group\_by\_dynamic}).

\begin{lstlisting}[style=py]
shop = sdata.get_example("shop")
shop.turn_into_dates("date")
shop.add_column_with_part_of_date("weekday", from_column="date", part="weekday name")
shop.add_column("check_on", from_column="date", plus="1 month")
shop.find_total_of("sales", for_each_group_in=["shop", "date"], every="month")
\end{lstlisting}

\begin{lstlisting}[style=py]
logger = sdata.get_example("fjord_logger")
logger.turn_into_dates("time_utc")
logger.change_time_zone_in("time_utc", to="Europe/Oslo", recorded_in="UTC")
logger.round_dates_in("time_utc", down_to="hour")
logger.show_first_rows(3)
\end{lstlisting}

\textbf{A date from year, month and day columns} (pandas \code{to\_datetime} of three columns, Polars
\code{pl.date}). Join the parts with a dash, then turn the text into dates:

\begin{lstlisting}[style=py]
table = sdata.table({"year": [2026, 2026], "month": [5, 6], "day": [12, 3]})
table.add_column("date", by_joining_text=["year", "month", "day"], separator="-")
table.turn_into_dates("date", written_like="2026-5-31")
table.show()
\end{lstlisting}

\textbf{Left to pandas and Polars.}
\begin{itemize}
  \item \textbf{Periods} such as ``2026-03'' as one month (pandas \code{to\_period}). Use
    \code{every="month"} in the summary words, or the part \code{"year and month"}.
  \item \textbf{Times of day on their own} (``08:30''), and rows at a time of day (pandas
    \code{between\_time}).
  \item \textbf{Working days} (NumPy \code{busday\_count}, Polars \code{business\_day\_count}).
  \item \textbf{Unix time}, seconds since 1970 (pandas \code{to\_datetime(unit="s")}).
\end{itemize}

\begin{lstlisting}[style=py]
import numpy as np
import pandas as pd

table = sdata.table({"start": ["2026-05-04", "2026-05-11"], "end": ["2026-05-15", "2026-05-22"],
                     "logged": [1778140800, 1778745600]})
df = table.as_pandas()
df["working_days"] = np.busday_count(df["start"].to_numpy("datetime64[D]"),
                                     df["end"].to_numpy("datetime64[D]"))
df["logged"] = pd.to_datetime(df["logged"], unit="s")
df["month"] = pd.to_datetime(df["start"]).dt.to_period("M").astype(str)
table = sdata.table(df)
table.show()
\end{lstlisting}

\unsection{Your own functions}{from-pandas:functions}

\textbf{The job.} Run a Python function that no word has. In the guides:
\href{https://pandas.pydata.org/docs/user_guide/user_defined_functions.html}{User-Defined Functions};
\href{https://docs.pola.rs/user-guide/expressions/user-defined-python-functions/}{User-defined Python
functions} in Polars.

\textbf{In simplydata.} \code{add\_column(..., by\_running=f)} runs your function on each row (pandas
\code{apply(axis=1)}, Polars \code{map\_elements}). The function gets the row as a dict.

\begin{lstlisting}[style=py]
def grade(row):
    return "top" if row["marks"] >= 80 else "other"

table = sdata.get_example("marks")
table.add_column("grade", by_running=grade)
table.count_each("grade")
\end{lstlisting}

\textbf{Left to pandas and Polars.} \code{pipe(f)}, a function on the whole table inside a chain.
simplydata words are already a chain of lines; a plain function that takes the table does the same:

\begin{lstlisting}[style=py]
def clean(table):
    table.clean_column_names()
    table.remove_rows_that_are_all_empty()
    return table

stations = clean(sdata.get_example("fjord_stations"))
stations.get_column_names()
\end{lstlisting}

\unsection{Charts and styled tables}{from-pandas:charts}

pandas draws charts with \code{df.plot}
(\href{https://pandas.pydata.org/docs/user_guide/visualization.html}{Chart visualization}). In
simplydata, \code{table.start\_chart()} gives a chart and its \code{draw\_} words; the charts part of
this manual has them all. Styled tables (pandas
\href{https://pandas.pydata.org/docs/user_guide/style.html}{Table Visualization}: colours, highlights,
number formats) are left to pandas; it needs the package jinja2. \code{table.save("table.html")} saves a
plain web page table.

\unsection{Speed, memory and copies}{from-pandas:speed}

\textbf{The job.} Work fast on big tables, keep memory low, and know when a table is a copy. In the
guides: \href{https://pandas.pydata.org/docs/user_guide/enhancingperf.html}{Enhancing performance},
\href{https://pandas.pydata.org/docs/user_guide/scale.html}{Scaling to large datasets},
\href{https://pandas.pydata.org/docs/user_guide/copy_on_write.html}{Copy-on-Write}; the
\href{https://docs.pola.rs/user-guide/concepts/lazy-api/}{Lazy API} in Polars.

\textbf{In simplydata.} \code{sdata.table("file", fast=True)} runs the table on Polars; the words and
the answers stay the same. A word changes its own table only: \code{backup = table.make\_copy()} keeps
the old one, and \code{table.undo()} goes back one step.

\textbf{Left to pandas and Polars.} Lazy plans and streaming (Polars \code{LazyFrame},
\code{collect(engine="streaming")}), memory use and smaller types (see
\hyperref[from-pandas:kinds]{Kinds of columns}), and sparse storage. Hand the table over for them:

\begin{lstlisting}[style=py]
import polars as pl

river = sdata.get_example("river")
plan = river.as_polars().lazy().filter(pl.col("nitrate") > 3).group_by("station").agg(pl.len())
result = sdata.table(plan.collect())
result.sort_by("station")
result.show()
\end{lstlisting}

\textbf{The numbers as a NumPy array} (pandas and Polars \code{to\_numpy}). Keep the number columns and
hand them over:

\begin{lstlisting}[style=py]
table = sdata.get_example("marks")
table.keep_columns(of_kind="number")
grid = table.as_pandas().to_numpy()
print(grid.shape)
\end{lstlisting}

\textbf{A correlation matrix} (pandas and Polars \code{corr}). This is a statistics word: name the
columns, and \code{find\_correlation\_between} gives every pair with its test.

\begin{lstlisting}[style=py]
table = sdata.get_example("marks")
table.keep_rows_where("subject", is_equal_to="Maths")      # one row for each child
table.find_correlation_between("hours", "marks", "height")
\end{lstlisting}

\unsection{SQL}{from-pandas:sql}

Polars runs SQL on its tables (\href{https://docs.pola.rs/user-guide/sql/intro/}{SQL}); pandas reads
SQL from a database with \code{read\_sql}. In simplydata, \code{table.run\_sql} runs a query on the
table, which is called \code{data} inside the query, and \code{sdata.run\_sql} runs one on files. DuckDB
does the work (pip install ``simplydata[sql]'').

% needs: duckdb
\begin{lstlisting}[style=py]
marks = sdata.get_example("marks")
best = marks.run_sql("SELECT class, count(*) AS children, max(marks) AS best FROM data "
                     "WHERE subject = 'Maths' GROUP BY class ORDER BY class")
best.show()
\end{lstlisting}

\unsection{The jobs left out, in one list}{from-pandas:left-out}

simplydata leaves these jobs to pandas and Polars on purpose. Each is shown above with the hand-over.
\begin{description}[style=nextline,leftmargin=1.2em]
  \item[The row index] set and reset an index, choose and slice rows by label, a MultiIndex, sort by the
    index, repeated labels, and lining up two tables by the index (\code{add(fill\_value=0)}).
    \hyperref[from-pandas:choosing]{Choosing rows and columns}, \hyperref[from-pandas:joining]{Joining}.
  \item[Functions in a chain] choose rows with a function (\code{loc[lambda d: ...]}), your own function on
    each group (\code{groupby().apply}), and \code{pipe}. \hyperref[from-pandas:groups]{Groups},
    \hyperref[from-pandas:functions]{Your own functions}.
  \item[Storage] memory use, smaller types, codes of categories, periods, lazy plans and streaming.
    \hyperref[from-pandas:kinds]{Kinds of columns}, \hyperref[from-pandas:categories]{Categories},
    \hyperref[from-pandas:dates]{Dates and times}, \hyperref[from-pandas:speed]{Speed}.
  \item[Other packages] the numbers as a NumPy array; the correlation matrix is a statistics word.
    \hyperref[from-pandas:speed]{Speed, memory and copies}.
\end{description}
The Polars guide has a page for readers who move the other way, from pandas to Polars:
\href{https://docs.pola.rs/user-guide/migration/pandas/}{Coming from pandas}.

\paragraph{How complete is simplydata?}
We typed every public name of these parts of the library into \code{sdata.help()} (the versions of November 2026).
pandas: DataFrame, Series, groupby, \code{.str}, \code{.dt}, the functions and classes at the top level, and the
rolling, expanding, ewm and resample objects, 800 names. 670 lead to a simplydata word or option, 93 hand over to
pandas with a note, and 37 are left to pandas on purpose because simplydata has no row index and no \code{pipe}.
Polars: DataFrame, Series, expressions, \code{.str}, \code{.dt}, \code{.list}, group\_by and the functions and
classes at the top level, 958 names. 787 lead to a word or option, 98 hand over to Polars with a note (such as its
exponentially weighted spread, folds, struct and JSON columns, engines and settings), and 73 are left to Polars on
purpose: list columns (simplydata keeps one value in each cell) and lazy plans. Every name has an answer in
\code{sdata.help()}.
